% ---------------------------------------------------------------------------
% Chapitre 09 — Les treillis simples : plats, ensembles bornés, intervalles
% Sources : notes/04-domains-lattices.md (§1–§3.5, §4.2–§4.4, §6.1–§6.3, §6.7,
% §8), src/domains/mod.rs, src/domains/impls/{mod,bool_val,setter_val,
% str_const,interval}.rs, src/domains/transfer/state_value.rs (eval_binop,
% eval_bitwise), src/engine/cfg_analyzer.rs (narrow_env_for_branch),
% src/engine/fixpoint.rs (seuils, boucle externe), ADR-002, ADR-014, ADR-020,
% commit 548f922.
% Sorties observées au commit e67b10a avec target/debug/reactant sur les
% fichiers /tmp/ch09/*.tsx.
% ---------------------------------------------------------------------------
\chapter{Les treillis simples : plats, ensembles bornés, intervalles}
\label{chap:treillis-simples}

\begin{objectifs}
  \item Lire le trait \code{AbstractDomain} comme un contrat : ce que chaque méthode doit garantir pour que l'analyse reste sound.
  \item Connaître les trois familles de treillis \og simples\fg{} du code : les treillis plats (\code{BoolVal}, \code{SetterVal}), l'ensemble borné de chaînes (\code{StrConst}), les intervalles (\code{Interval}).
  \item Savoir calculer à la main toutes les opérations des intervalles : ordre, join, arithmétique (\texttt{+ - * \% **}), opérateurs bit à bit, widening simple et à seuils, raffinement sur gardes.
  \item Savoir prouver en quelques lignes la soundness de chacune de ces opérations, et reconnaître les endroits où le code s'en écarte.
  \item Connaître les limites observées : $\code{NaN}$ et $\pm\infty$, $0 \times \pm\infty$, décalages sur des non-entiers, cycles finis.
\end{objectifs}

\prerequis{\cref{chap:interpretation-abstraite} (treillis, concrétisation, widening, soundness : \cref{def:treillis}, \cref{def:concretisation}, \cref{def:widening}, \cref{def:soundness}) ; \cref{chap:ir} pour les expressions \code{Expr}, \code{BinOp}, \code{Prim} ; \cref{chap:lowering-cfg} pour la forme des conditions de branche.}

Le \cref{chap:interpretation-abstraite} a présenté les treillis \og en un coup
d'œil\fg{} et suivi un compteur jusqu'à son point fixe. Ce chapitre ouvre la
boîte. Il décrit, opération par opération, les quatre domaines de base sur
lesquels repose toute valeur abstraite de \reactant{}. Ces domaines sont
petits (moins de 900 lignes à eux quatre), mais chaque ligne compte : une
borne mal arrondie ou un cas $\Bot$ oublié suffit à produire un faux négatif,
et un faux négatif est interdit (\cref{def:soundness}).

Le \cref{chap:statevalue-stabilite} assemblera ces domaines en un produit,
\code{StateValue}, et y ajoutera le treillis de stabilité référentielle. On
peut déjà situer les pièces (\cref{fig:treillis-simples-produit}) : sur les
huit composantes de \code{StateValue}, quatre sont les treillis de ce
chapitre, trois sont de simples booléens \og cette sorte est-elle possible
?\fg{}, et une est la stabilité.

\begin{figure}[htbp]\centering
\begin{tikzpicture}[node distance=2mm]
  \node[bloc,minimum width=26mm] (num) {\code{num}\\\code{Interval}};
  \node[bloc,minimum width=26mm,right=of num] (bool) {\code{boolean}\\\code{BoolVal}};
  \node[bloc,minimum width=26mm,right=of bool] (str) {\code{str}\\\code{StrConst}};
  \node[bloc,minimum width=26mm,right=of str] (set) {\code{setter}\\\code{SetterVal}};
  \node[cfgbloc,minimum width=26mm,below=4mm of num] (ref) {\code{reference}\\\code{Stability}};
  \node[cfgbloc,minimum width=26mm,right=of ref] (null) {\code{null}\\\code{bool}};
  \node[cfgbloc,minimum width=26mm,right=of null] (undef) {\code{undef}\\\code{bool}};
  \node[cfgbloc,minimum width=26mm,right=of undef] (other) {\code{other}\\\code{bool}};
  \node[draw=ra-gray,dashed,rounded corners,fit=(num)(set)(ref)(other),inner sep=4pt,label={[etiquette]above:\code{StateValue} (\src{src/domains/impls/state_value.rs}{26}{44})}] {};
\end{tikzpicture}
\caption{Les huit composantes du produit \code{StateValue}. En bleu, les
treillis étudiés dans ce chapitre ; en gris, ceux du
\cref{chap:statevalue-stabilite} (la stabilité et les trois bits de sorte).}
\label{fig:treillis-simples-produit}
\end{figure}

% ===========================================================================
\section{Le contrat d'un domaine : le trait \code{AbstractDomain}}
\label{sec:treillis-simples-contrat}

\subsection{Une question concrète}

Prenons le programme qui servira de fil rouge à tout le chapitre.

\begin{lstlisting}[language=TypeScript,caption={Deux compteurs (\file{/tmp/ch09/counter.tsx}, repris de \file{tests/fixtures/widening.tsx})},label={lst:treillis-simples-counter}]
import { useState, useEffect } from "react";

export function Unbounded() {
  const [count, setCount] = useState(0);
  useEffect(() => {
    setCount(count + 1);
  }, [count]);
  return <div>{count}</div>;
}

export function Guarded() {
  const [count, setCount] = useState(0);
  useEffect(() => {
    if (count < 10) setCount(count + 1);
  }, [count]);
  return <div>{count}</div>;
}
\end{lstlisting}

\code{Unbounded} boucle : chaque exécution de l'effet écrit une nouvelle valeur
dans \code{count}, ce qui provoque un nouveau rendu, donc une nouvelle
exécution de l'effet. \code{Guarded} s'arrête à 10. Pour les distinguer,
l'analyseur doit répondre à une question sur des \emph{nombres} : l'ensemble
des valeurs que \code{count} peut prendre est-il borné ? Il lui faut donc un
domaine qui représente des ensembles de nombres (les intervalles), qui sache
calculer \code{count + 1} (l'arithmétique), exploiter \code{count < 10} (le
raffinement sur garde) et converger malgré une suite de valeurs infinie (le
widening). L'analyseur répond ainsi :

\begin{sortie}
$ reactant check /tmp/ch09/counter.tsx --show-clean
  Guarded  (2 hooks)  /tmp/ch09/counter.tsx  ✓
  Unbounded  (2 hooks)  /tmp/ch09/counter.tsx
    warn   infinite-loop  [hook:0]  (line 5:2)  this effect keeps pushing state `count` (its deps do not provably gate it, so the effect can re-run every render) to new values on every run. Potential infinite render loop
       (2 trace step(s), rerun with --trace)

⚠  1 warning(s) across 1 file(s).
\end{sortie}

Toutes les pièces de cette réponse sont dans ce chapitre, sauf la règle
elle-même (\cref{chap:infinite-loop}) et la boucle de point fixe
(\cref{chap:point-fixe-composant}).

\subsection{Le trait, méthode par méthode}

Un domaine abstrait est, dans le code, un type qui implémente
\code{AbstractDomain} (\file{src/domains/mod.rs}). La tête du trait a été
présentée au \cref{chap:interpretation-abstraite}
(\cref{lst:interpretation-abstraite-abstractdomain}) : \code{bottom},
\code{top}, \code{is_bottom}, \code{join}, \code{meet}, \code{widen}, et
\code{widen_to} dont l'implémentation par défaut ignore les seuils. L'ordre du
treillis n'est pas une méthode : c'est le supertrait \code{PartialOrd}
(\code{a <= b} signifie $a \lleq b$, \code{partial_cmp} rend \code{None} pour
deux éléments incomparables). La suite du trait est moins connue ; la voici.

\rustsnippet[label={lst:treillis-simples-narrow-defaut}]{src/domains/mod.rs}{40}{85}{ponts vers \code{StateValue} et raffinements par défaut}

Suivent \code{narrow_leq}, \code{narrow_gt}, \code{narrow_geq},
\code{narrow_eq} et \code{narrow_neq}, de la même forme
(\src{src/domains/mod.rs}{86}{100}). Trois groupes de méthodes :
\begin{itemize}
  \item \textbf{Les ponts} \code{as_state_value} / \code{from_state_value}
    servent au tas, qui range l'environnement capturé par une fermeture sous
    forme de \code{StateValue}. Seul \code{StateValue} les redéfinit.
  \item \textbf{Les raffinements de sorte} (\code{narrow_drop_null},
    \code{narrow_drop_undef}, \code{narrow_keep_nullish_only},
    \code{narrow_truthy}, \code{narrow_falsy}) ne concernent que le produit,
    seul à savoir ce qu'est une sorte ; ils sont étudiés au
    \cref{chap:statevalue-stabilite}, sauf \code{narrow_truthy} dont l'effet
    sur les nombres est vu en \cref{sec:treillis-simples-gardes}.
  \item \textbf{Les raffinements numériques} \code{narrow_lt} \dots{}
    \code{narrow_neq} restreignent une valeur à celles qui satisfont
    \code{x < v}, \code{x <= v}, etc., pour un littéral \code{v}.
    \code{Interval} les redéfinit, et \code{StateValue} les délègue à sa
    composante \code{num}.
\end{itemize}

Tous les raffinements valent l'identité par défaut, et le commentaire dit
pourquoi c'est permis : ne rien raffiner est toujours sound, seulement
imprécis. Le \cref{tab:treillis-simples-qui-redefinit} récapitule qui redéfinit
quoi.

\begin{table}[htbp]\centering\small
\begin{tabular}{lcccc}
\hline
Domaine & \code{widen_to} & ponts & \code{narrow_lt} \dots{} \code{narrow_neq} & raffinements de sorte \\
\hline
\code{BoolVal}, \code{SetterVal} (macro) & défaut & défaut & identité & identité \\
\code{StrConst} & défaut & défaut & identité & identité \\
\code{Stability} & défaut & défaut & identité & identité \\
\code{Interval} & redéfini & défaut & redéfinis & identité \\
\code{StateValue} & redéfini & redéfinis & délégués à \code{num} & redéfinis \\
\hline
\end{tabular}
\caption{Les méthodes à implémentation par défaut de \code{AbstractDomain}, et
les domaines qui les redéfinissent (relevé des \code{impl AbstractDomain}).}
\label{tab:treillis-simples-qui-redefinit}
\end{table}

\subsection{Le contrat}

Le trait ne dit rien des propriétés que ses méthodes doivent avoir ; elles
sont pourtant le fondement de la soundness. Écrivons-les une fois.

\begin{definition}{Contrat d'un domaine abstrait}{treillis-simples-contrat}
  Soit $D$ un type qui implémente \code{AbstractDomain}, muni d'une
  concrétisation $\gam : D \to \powerset(\Val)$ (\cref{def:concretisation}),
  où $\Val$ est l'ensemble des valeurs JavaScript concernées. Le domaine
  respecte le contrat si, pour tous $a, b \in D$ :
  \begin{enumerate}
    \item \emph{ordre} : $a \lleq b \Rightarrow \gam(a) \subseteq \gam(b)$ ;
    \item \emph{bottom et top} : $\gam(\Bot) = \emptyset$ et $\gam(\Top) = \Val$ ;
      \code{is_bottom(a)} n'est vrai que si $\gam(a) = \emptyset$ ;
    \item \emph{join} : $a \lleq a \join b$ et $b \lleq a \join b$, donc
      $\gam(a) \cup \gam(b) \subseteq \gam(a \join b)$ ;
    \item \emph{widening} : \code{widen} et \code{widen_to} (pour tout
      ensemble fini de seuils) majorent le join et rendent toute suite
      d'élargissements stationnaire (\cref{def:widening}) ;
    \item \emph{raffinement} : pour une condition $c$ et son opérateur
      \code{narrow_c}, $\gam(\code{narrow_c}(a)) \supseteq \gam(a) \cap
      \{x \mid c(x)\}$ ;
    \item \emph{meet} : $\gam(a) \cap \gam(b) \subseteq \gam(a \meet b)$.
  \end{enumerate}
\end{definition}

Les conditions 3 à 5 portent la soundness : un join ou un raffinement qui
perdrait une valeur concrète ferait disparaître un comportement réel. Le join
n'a pas besoin d'être la plus petite borne supérieure pour être sound (il suffit
qu'il majore), mais il l'est pour les quatre domaines du chapitre, et c'est ce
qui fait leur précision. La condition 6 est de pure forme aujourd'hui :

\begin{piege}
  \code{meet} n'est appelé nulle part hors des tests et des délégations
  internes (\code{StateValue::meet} vers ses composantes). Le raffinement de
  l'analyse passe exclusivement par les \code{narrow_*}, jamais par
  \code{meet}. Un contributeur qui voudrait raffiner un état par
  \og intersection avec la garde\fg{} doit donc écrire un \code{narrow_*}, ou
  vérifier d'abord que le \code{meet} du domaine concerné respecte la
  condition 6 (ce n'est pas le cas de \code{Stability},
  \cref{chap:statevalue-stabilite}).
\end{piege}

\begin{piege}
  Dans \code{impl AbstractDomain for Interval}, les raffinements du trait sont
  écrits \code{Interval::narrow_lt(&self, v)} et non \code{self.narrow_lt(v)}
  (\src{src/domains/impls/interval.rs}{362}{365}). La méthode du trait prend
  \code{self} par valeur, la méthode inhérente \code{&self}. Dans le corps du
  trait, \code{self.narrow_lt(v)} serait résolu dès l'étape \og par valeur\fg{}
  de la recherche de méthode, donc vers la méthode du trait elle-même :
  récursion infinie. Le commentaire \og fully-qualified\fg{} du code le
  signale.
\end{piege}

\begin{remarque}
  Le commentaire en tête du trait annonce encore \code{Clone + Copy}, alors
  que la borne réelle ne contient pas \code{Copy}
  (\cref{chap:interpretation-abstraite}). Des quatre domaines de ce chapitre,
  \code{BoolVal} et \code{Interval} sont \code{Copy} ; \code{SetterVal} et
  \code{StrConst} ne dérivent pas \code{Copy} (le second contient un
  \code{Arc}), d'où les \code{clone()} de la macro ci-dessous.
\end{remarque}

% ===========================================================================
\section{Les treillis plats : \code{BoolVal} et \code{SetterVal}}
\label{sec:treillis-simples-plats}

\subsection{Deux exemples React}

Un booléen d'état est l'exemple le plus simple de valeur à suivre :

\begin{lstlisting}[language=TypeScript,caption={Un booléen d'état et un setter choisi dynamiquement (exemple d'illustration)},label={lst:treillis-simples-bool}]
export function Panel({ admin }: { admin: boolean }) {
  const [open, setOpen] = useState(false);
  const [a, setA] = useState(0);
  const [b, setB] = useState(0);
  const set = admin ? setA : setB;
  return <button onClick={() => { setOpen(!open); set(1); }}>{open ? "-" : "+"}</button>;
}
\end{lstlisting}

Que sait-on de \code{open} ? Au premier rendu, qu'il vaut \code{false} ;
après un clic, qu'il vaut \code{true} ; en général, qu'il vaut l'un ou
l'autre. Que sait-on de \code{set} ? Que c'est \emph{un} setter de
\code{useState}, mais lequel dépend de la prop : \code{setA} ou \code{setB}.
Dans les deux cas, l'information utile est \og une valeur précise, ou bien on
ne sait pas laquelle\fg{}. C'est exactement ce que capture un treillis plat.

\begin{definition}{Treillis plat}{treillis-simples-plat}
  Soit $M$ un ensemble (les éléments \og intermédiaires\fg{}). Le
  \terme{treillis plat} sur $M$ est $M \cup \{\Bot, \Top\}$ ordonné par
  $\Bot \lleq m \lleq \Top$ pour tout $m \in M$, deux éléments distincts de
  $M$ étant incomparables. Son join vaut $a \join a = a$,
  $\Bot \join x = x \join \Bot = x$, et $\Top$ dans tous les autres cas ; son
  meet est le dual. Sa hauteur est 2 (plus longue chaîne stricte :
  $\Bot \sqsubset m \sqsubset \Top$).
\end{definition}

\Cref{fig:treillis-simples-hasse-plat} dessine le treillis plat générique et
indique, sur chaque arête, quel bras du code la justifie.

\begin{figure}[htbp]\centering
\begin{tikzpicture}[treillis]
  \node (bot) at (0,0) {$\Bot$};
  \node (m1) at (-2.2,1.6) {$m_1$};
  \node (m2) at (-0.8,1.6) {$m_2$};
  \node (m3) at (0.6,1.6) {$m_3$};
  \node (md) at (1.8,1.6) {$\cdots$};
  \node (top) at (0,3.2) {$\Top$};
  \foreach \n in {m1,m2,m3} { \draw (bot) -- (\n) -- (top); }
  \draw[dashed,ra-gray] (bot) -- (md) -- (top);
  \node[etiquette,align=left,anchor=west] at (2.4,2.6) {$m_i \lleq \Top$ : bras \code{(_, $top)}};
  \node[etiquette,align=left,anchor=west] at (2.4,0.6) {$\Bot \lleq m_i$ : bras \code{($bottom, _)}};
  \node[etiquette,align=left,anchor=west] at (2.4,1.6) {$m_i \neq m_j$ : aucun bras, \code{None}};
  \begin{scope}[xshift=-5.2cm]
    \node (bb) at (0,0) {\code{Bottom}};
    \node (bt) at (-1.1,1.6) {\code{True}};
    \node (bf) at (1.1,1.6) {\code{False}};
    \node (btop) at (0,3.2) {\code{Top}};
    \draw (bb) -- (bt) -- (btop) -- (bf) -- (bb);
    \node[etiquette] at (0,-0.5) {$\gam = \emptyset$};
    \node[etiquette] at (0,3.7) {$\gam = \{\mathtt{true},\mathtt{false}\}$};
  \end{scope}
\end{tikzpicture}
\caption{À gauche, \code{BoolVal}, dont la concrétisation est une bijection
avec $\powerset(\{\mathtt{true},\mathtt{false}\})$. À droite, le treillis plat
engendré par \code{flat_lattice!} pour un ensemble quelconque d'éléments
intermédiaires, avec le bras de \code{partial_cmp} qui établit chaque
relation.}
\label{fig:treillis-simples-hasse-plat}
\end{figure}

\subsection{Ce que fait le code : la macro \code{flat_lattice!}}

Les deux treillis plats du code ne diffèrent que par la charge utile de leurs
éléments intermédiaires. Leur ordre et leurs opérations sont donc engendrés
par une seule macro, introduite lors de la campagne de dette technique
(décision D5 d'\adr{20}).

\rustsnippet[label={lst:treillis-simples-flat}]{src/domains/impls/mod.rs}{1}{53}{la macro \code{flat_lattice!}}

Lisons-la bras par bras.
\begin{itemize}
  \item \textbf{Paramètres} (l.~11). La macro reçoit le type et les noms de
    ses variantes $\Bot$ et $\Top$ ; toutes les autres variantes sont les
    éléments intermédiaires. Elle exige \code{PartialEq + Clone} (les gardes
    \code{a == b} et les \code{clone()}).
  \item \textbf{\code{partial_cmp}} (l.~13--21). L'égalité d'abord ; puis
    \og $\Bot$ à gauche ou $\Top$ à droite\fg{} donne \code{Less} ; puis
    \og $\Top$ à gauche ou $\Bot$ à droite\fg{} donne \code{Greater} ; le reste,
    deux intermédiaires distincts, est incomparable. L'ordre des bras compte
    pour le couple $(\Bot, \Top)$, qui satisfait les deux motifs : le premier
    gagne, et c'est bien \code{Less}.
  \item \textbf{\code{join}} (l.~34--40). Égaux : on garde. L'un est $\Bot$ :
    on garde l'autre (le motif \code{($bottom, x) | (x, $bottom)} lie
    \code{x} des deux côtés). Sinon : $\Top$. Ce \og sinon\fg{} couvre aussi les
    cas où l'un est $\Top$, ce qui est correct.
  \item \textbf{\code{meet}} (l.~41--47) est le dual : $\Top$ est neutre, deux
    éléments différents donnent $\Bot$.
  \item \textbf{\code{widen}} (l.~48--50) est le join : la hauteur est finie.
\end{itemize}

\begin{proposition}{\code{flat_lattice!} engendre un treillis}{treillis-simples-flat-treillis}
  Pour tout type instancié par \code{flat_lattice!}, \code{partial_cmp} est
  l'ordre du treillis plat de la \cref{def:treillis-simples-plat}, \code{join}
  en est la plus petite borne supérieure, \code{meet} la plus grande borne
  inférieure, et \code{widen} est un widening.
\end{proposition}
\begin{proof}
  L'ordre : les trois bras reproduisent exactement les trois clauses de la
  définition (réflexivité, $\Bot$ minimum, $\Top$ maximum), le dernier bras
  déclare incomparables deux intermédiaires distincts. Le join : si
  $a = b$, $a$ est la plus petite borne de $\{a\}$ ; si $a = \Bot$, $b$ majore
  $a$ et $b$ et tout majorant de $b$ en est un de $\{a,b\}$ ; sinon, soit l'un
  vaut $\Top$ (seul majorant), soit $a$ et $b$ sont deux intermédiaires
  distincts, dont le seul majorant commun est $\Top$. Le meet est dual. Le
  widening : \code{widen = join} majore le join, et toute suite croissante
  d'un treillis de hauteur 2 est stationnaire après au plus deux pas stricts.
\end{proof}

La macro est définie \emph{avant} les déclarations \code{pub mod bool_val;} et
\code{pub mod setter_val;} (\src{src/domains/impls/mod.rs}{55}{60}). C'est la
portée \emph{textuelle} des \code{macro_rules!} qui la rend visible dans ces
sous-modules, sans \texttt{\#[macro\_export]} ni \code{use}. Déplacer la
déclaration des modules au-dessus de la macro casserait la compilation.

\subsection{\code{BoolVal} : un treillis plat exact}

\rustsnippet[label={lst:treillis-simples-boolval}]{src/domains/impls/bool_val.rs}{1}{13}{le treillis des booléens}

La concrétisation est
$\gam(\code{Bottom}) = \emptyset$, $\gam(\code{True}) = \{\mathtt{true}\}$,
$\gam(\code{False}) = \{\mathtt{false}\}$,
$\gam(\code{Top}) = \{\mathtt{true}, \mathtt{false}\}$. C'est une bijection
avec les parties de $\{\mathtt{true}, \mathtt{false}\}$, compatible avec
l'ordre : le treillis plat \emph{est} ici le treillis des parties, et join et
meet sont exacts ($\gam(a \join b) = \gam(a) \cup \gam(b)$). C'est le seul
domaine du chapitre sans aucune perte.

Où naissent les booléens abstraits ? Un littéral \code{true} ou \code{false}
donne \code{True} ou \code{False}, comme un littéral numérique donne
\code{Interval::point} et une chaîne un singleton : trois cas d'un même
\code{match} (\src{src/domains/transfer/state_value.rs}{112}{117}).

La négation \code{!x} inverse exactement un
booléen pur (\src{src/domains/transfer/state_value.rs}{947}{965}). En
revanche, \emph{toute} comparaison (\code{==}, \code{<}, \code{in},
\code{instanceof}\dots) rend \code{BoolVal::Top}, même \code{1 < 2}
(\src{src/domains/transfer/state_value.rs}{797}{806}) : le domaine ne
calcule pas la valeur d'une comparaison ; c'est le raffinement de branche
(\cref{sec:treillis-simples-gardes}) qui exploite la condition. Rendre un
booléen pur, et non $\Top$ tout entier, a son importance : les composantes
\code{num} et \code{str} du résultat restent à $\Bot$, ce qui permet
d'enchaîner d'autres raffinements sur le résultat.

\begin{piege}
  \code{!x} sur un $x$ qui vaut $\Bot$ rend \code{boolean(Top)} et non
  $\Bot$ : le test \og booléen pur\fg{} de la ligne 949 est satisfait par
  $\Bot$ (toutes ses composantes sont vides), puis le \code{match} tombe dans
  \code{_ => BoolVal::Top}. Un chemin mort redevient vivant pour la suite
  de l'évaluation. C'est imprécis, pas unsound : $\Top$ majore tout.
\end{piege}

\subsection{\code{SetterVal} : un treillis plat à identité}

\rustsnippet[label={lst:treillis-simples-setterval}]{src/domains/impls/setter_val.rs}{3}{37}{le treillis plat des setters}

\begin{itemize}
  \item \code{One(ComponentId, HookLabel)} désigne \emph{le} setter du slot
    \code{HookLabel} (\cref{def:slot}) du composant \code{ComponentId} :
    $\gam(\code{One}(c,l)) = \{\text{la fonction }\code{setState}\text{ du
    slot } l \text{ de } c\}$. L'identité est un \code{ComponentId} interné
    (\adr{40}), plus un nom de composant comme dans les premières versions.
  \item \code{Top} : $\gam$ est l'ensemble de tous les setters de
    \code{useState}. On sait que la valeur est un setter, pas lequel.
  \item \code{as_one} (l.~27--32) est l'accesseur qu'utilisent les règles et
    les relations pour savoir \emph{quel} slot un appel écrit : il ne répond
    que si l'identité est exacte.
\end{itemize}

Dans le \cref{lst:treillis-simples-bool}, les deux branches de
\code{admin ? setA : setB} donnent \code{One(c,1)} et \code{One(c,2)} ; leur
join est \code{Top}. Ici le treillis plat \emph{perd} de l'information :
$\gam(\code{Top})$ contient tous les setters, alors que la vérité est
\og \code{setA} ou \code{setB}\fg{}. La perte est sound (on majore) et c'est le
prix de la hauteur 2.

\begin{react}[identité des setters]
  React garantit que la fonction \code{setState} rendue par \code{useState}
  garde la même identité d'un rendu à l'autre. C'est pourquoi tout
  \code{SetterVal} non $\Bot$ est projeté sur la stabilité \code{Stable}
  (commentaire des lignes 13--14 et 35--36) : on peut ignorer quel setter,
  on sait qu'il ne change pas.
\end{react}

% ===========================================================================
\section{Un ensemble borné de chaînes : \code{StrConst}}
\label{sec:treillis-simples-strconst}

\subsection{Le problème}

Les chaînes d'état sont souvent des \emph{énumérations} déguisées :

\begin{lstlisting}[language=TypeScript,caption={Un mode d'affichage (exemple d'illustration)},label={lst:treillis-simples-mode}]
export function Gallery() {
  const [mode, setMode] = useState("grid");
  return (
    <div className={"gallery-" + mode}>
      <button onClick={() => setMode("list")}>list</button>
      <button onClick={() => setMode("grid")}>grid</button>
    </div>
  );
}
\end{lstlisting}

Un treillis plat sur les chaînes dirait de \code{mode} : $\Top$, puisque deux
constantes différentes y sont écrites. On aimerait garder
$\{\code{"grid"}, \code{"list"}\}$, et pour \code{className},
$\{\code{"gallery-grid"}, \code{"gallery-list"}\}$. Le treillis des parties de
l'ensemble des chaînes le permettrait, mais il est de hauteur infinie : une
chaîne qui grandit (\code{s + "x"}) engendrerait une suite d'ensembles
strictement croissante. La solution classique borne la taille des ensembles.

\begin{definition}{Ensemble $k$-borné}{treillis-simples-k-ensemble}
  Soit $S$ un ensemble (ici les chaînes) et $k \ge 1$. Le treillis des
  \terme{ensembles $k$-bornés} sur $S$ est
  $\{\Bot\} \cup \{X \subseteq S \mid 1 \le |X| \le k\} \cup \{\Top\}$,
  ordonné par l'inclusion, $\Bot$ en dessous de tout et $\Top$ au-dessus. La
  concrétisation est $\gam(\Bot) = \emptyset$, $\gam(X) = X$,
  $\gam(\Top) = S$. Le join de $X$ et $Y$ est $X \cup Y$ si
  $|X \cup Y| \le k$, et $\Top$ sinon ; le meet est $X \cap Y$ (et $\Bot$ si
  l'intersection est vide). La hauteur est $k + 1$.
\end{definition}

\subsection{Ce que fait le code}

\rustsnippet[label={lst:treillis-simples-strconst}]{src/domains/impls/str_const.rs}{7}{39}{\code{StrConst}, ses constructeurs, $k = 4$}

\begin{itemize}
  \item \code{STR_WIDEN_THRESHOLD} vaut 4 : c'est le $k$ de la définition.
    Ce seuil vient du domaine produit d'origine (\adr{8}), qui l'a transmis
    tel quel à \adr{15}.
  \item \texttt{Set(Arc<BTreeSet<String>{}>)} : l'ensemble est trié
    (\code{BTreeSet}), ce qui rend l'égalité et l'affichage déterministes, et
    partagé (\code{Arc}), ce qui rend le \code{clone} d'une valeur abstraite
    bon marché : les environnements sont copiés à chaque bloc.
  \item \code{singleton} fabrique $\{s\}$ ; \code{from_set} est le
    constructeur \emph{canonique} : vide $\mapsto \Bot$, plus de 4 éléments
    $\mapsto \Top$, sinon l'ensemble.
\end{itemize}

\begin{invariant}{Forme canonique de \code{StrConst}}{treillis-simples-strconst}
  Toute valeur \code{Set(s)} construite par \code{singleton} ou
  \code{from_set} vérifie $1 \le |s| \le 4$. L'invariant n'est pas imposé par
  le type : la variante étant publique, \code{StrConst::Set(Arc::new(}$\emptyset$\code{))}
  reste constructible, et \code{is_bottom} y répondrait \code{false}. Tout
  code qui fabrique un ensemble doit passer par \code{from_set}.
\end{invariant}

L'ordre et les opérations :

\rustsnippet[label={lst:treillis-simples-strconst-ord}]{src/domains/impls/str_const.rs}{42}{62}{l'ordre de \code{StrConst} : l'inclusion}

\rustsnippet[label={lst:treillis-simples-strconst-join}]{src/domains/impls/str_const.rs}{75}{102}{join, meet et widen de \code{StrConst}}

L'ordre est l'inclusion, avec les deux sentinelles aux extrémités ; les bras
$\Bot$ précèdent les bras $\Top$, si bien que $(\Bot, \Top)$ donne
\code{Less}. Le join calcule l'union puis la fait passer par \code{from_set},
qui la remplace par $\Top$ si elle a plus de 4 éléments. Le meet calcule
l'intersection, qui ne peut pas grossir, et que \code{from_set} transforme en
$\Bot$ si elle est vide.

\Cref{fig:treillis-simples-hasse-strconst} en dessine un fragment.

\begin{figure}[htbp]\centering
\begin{tikzpicture}[treillis,xscale=1.05]
  \node (bot) at (0,0) {$\Bot$};
  \node (a) at (-3,1.2) {$\{a\}$};
  \node (b) at (-1,1.2) {$\{b\}$};
  \node (c) at (1,1.2) {$\{c\}$};
  \node (d) at (3,1.2) {$\{d\}$};
  \node (ab) at (-3,2.4) {$\{a,b\}$};
  \node (bc) at (-1,2.4) {$\{b,c\}$};
  \node (cd) at (1,2.4) {$\{c,d\}$};
  \node (dots2) at (3,2.4) {$\cdots$};
  \node (abc) at (-2,3.6) {$\{a,b,c\}$};
  \node (bcd) at (0.5,3.6) {$\{b,c,d\}$};
  \node (dots3) at (2.6,3.6) {$\cdots$};
  \node (abcd) at (-0.5,4.8) {$\{a,b,c,d\}$};
  \node (abce) at (2.2,4.8) {$\{a,b,c,e\}$};
  \node (dots4) at (3.9,4.8) {$\cdots$};
  \node (top) at (0.6,6.0) {$\Top$};
  \foreach \n in {a,b,c,d} { \draw (bot) -- (\n); }
  \draw (a) -- (ab) (b) -- (ab) (b) -- (bc) (c) -- (bc) (c) -- (cd) (d) -- (cd);
  \draw (ab) -- (abc) (bc) -- (abc) (bc) -- (bcd) (cd) -- (bcd);
  \draw (abc) -- (abcd) (bcd) -- (abcd) (abc) -- (abce);
  \draw (abcd) -- (top) (abce) -- (top);
  \draw[dashed,ra-gray] (dots4) -- (top);
  \node[etiquette,anchor=west] at (4.3,6.0) {$\gam = $ toutes les chaînes};
  \node[etiquette,anchor=west] at (4.3,4.8) {taille 4 ($k$)};
  \node[etiquette,anchor=west] at (4.3,1.2) {singletons};
  \node[etiquette,anchor=west] at (4.3,0) {$\gam = \emptyset$};
  \draw[fleche,ra-amber] (abc.north west) to[bend left=35] node[etiquette,left,pos=0.6] {$\join \{d,e\}$} (top.west);
\end{tikzpicture}
\caption{Un fragment de \code{StrConst} (quelques arêtes seulement). Les
ensembles de taille 1 à 4 sont ordonnés par inclusion ; au-delà, le join
saute à $\Top$ : $\{a,b,c\} \join \{d,e\} = \Top$ (flèche), car l'union a cinq
éléments.}
\label{fig:treillis-simples-hasse-strconst}
\end{figure}

\begin{proposition}{\code{StrConst} est un treillis de hauteur 5}{treillis-simples-strconst}
  \code{partial_cmp} est l'ordre des ensembles 4-bornés, \code{join} en est
  la plus petite borne supérieure, \code{meet} la plus grande borne
  inférieure. La plus longue chaîne stricte a six éléments
  ($\Bot \sqsubset \{a\} \sqsubset \{a,b\} \sqsubset \{a,b,c\} \sqsubset
  \{a,b,c,d\} \sqsubset \Top$), donc \code{widen = join} est un widening.
\end{proposition}
\begin{proof}
  Pour le join de $X$ et $Y$ : tout majorant commun est soit $\Top$, soit un
  ensemble $Z$ de taille au plus 4 contenant $X$ et $Y$, donc contenant
  $X \cup Y$. Si $|X \cup Y| \le 4$, $X \cup Y$ est lui-même un élément et
  c'est le plus petit majorant. Sinon aucun ensemble représentable ne contient
  $X \cup Y$, et le seul majorant est $\Top$. Pour le meet, tout minorant
  commun est $\Bot$ ou un ensemble inclus dans $X$ et dans $Y$, donc dans
  $X \cap Y$, qui a au plus 4 éléments : c'est le plus grand minorant (ou
  $\Bot$ s'il est vide). La chaîne gagne au plus un élément par pas, et
  $\Top$ est atteint au cinquième ajout.
\end{proof}

La soundness du join découle de la proposition : il majore les deux
opérandes, donc $\gam(X) \cup \gam(Y) \subseteq \gam(X \join Y)$. Mais
l'égalité ne vaut plus au-delà du seuil :
$\gam(\{a,b,c\} \join \{d,e\}) = S \supsetneq \{a,b,c,d,e\}$. Le
\og widening au seuil\fg{} du commentaire du code n'est donc pas un widening
au sens de la \cref{def:widening} (c'est le join, qui suffit déjà), mais une
perte de précision inscrite dans le join lui-même.

\begin{piege}
  \code{StrConst} dérive \code{PartialEq}, et l'égalité d'un
  \texttt{Arc<BTreeSet<String>{}>} compare le \emph{contenu}, pas les pointeurs
  (l'\code{impl PartialEq for Arc<T>} délègue à \code{T}). Deux ensembles
  égaux alloués séparément sont égaux, ce qu'exige le test de convergence du
  point fixe. Remplacer cette égalité par \code{Arc::ptr_eq}
  \og pour aller plus vite\fg{} empêcherait l'analyse de converger.
\end{piege}

\subsection{Où naissent les ensembles de chaînes}

Un littéral chaîne donne un singleton. L'opérateur \code{+} sur deux chaînes
pures calcule le produit cartésien des concaténations, puis le passe par
\code{from_set} (\src{src/domains/transfer/state_value.rs}{754}{768}, qui
est le début de la table des opérateurs du
\cref{lst:treillis-simples-eval-binop}). Dans le
\cref{lst:treillis-simples-mode}, \code{"gallery-" + mode} avec
$\code{mode} = \{\code{"grid"}, \code{"list"}\}$ donne
$\{\code{"gallery-grid"}, \code{"gallery-list"}\}$. Si l'un des opérandes
est une chaîne inconnue ($\Top$), le résultat est \og une chaîne\fg{}
(\code{StateValue::str_top()}), et non $\Top$ tout entier. Enfin
\code{typeof x} rend un singleton exact (\code{"number"},
\code{"string"}\dots) dès que $x$ n'a qu'une sorte possible
(\src{src/domains/transfer/state_value.rs}{969}{977}).

\begin{decision}[ADR-20 --- pas de combinateur \code{BoundedPowerset}]
  La campagne de dette technique a examiné l'extraction d'un combinateur
  générique \code{BoundedPowerset<T, N>} dont \code{StrConst} serait une
  instance (non-changement n\textsuperscript{o}~5 d'\adr{20}). Refusé : seul
  \code{StrConst} est un pur ensemble borné. \code{Stability}, qui ressemble à
  un ensemble borné de slots, mêle à ce fragment des éléments
  (\code{PerRender}, \code{Unknown}) aux joins croisés ; la faire entrer dans
  le combinateur restructurerait un type critique pour la soundness pour un
  seul consommateur réel. Abstraction prématurée.
\end{decision}

% ===========================================================================
\section{Les intervalles : représentation, ordre, join}
\label{sec:treillis-simples-intervalles}

\subsection{Le problème}

Revenons à \code{Guarded} (\cref{lst:treillis-simples-counter}). L'ensemble
exact des valeurs de \code{count} est $\{0, 1, \dots, 10\}$. Pour
\code{Unbounded}, c'est $\Nat$ tout entier. Un domaine de parties finies ne
peut représenter le second, et le premier serait déjà coûteux avec une borne
à $10^6$. Les \emph{intervalles} gardent seulement la plus petite et la plus
grande valeur : $[0, 10]$ et $[0, +\infty]$. On perd les trous (un intervalle
est \emph{convexe}), mais on garde exactement ce que la règle
\regle{infinite-loop} veut savoir : la valeur est-elle bornée ?

\begin{definition}{Intervalle abstrait}{treillis-simples-intervalle}
  Un intervalle abstrait est un triplet $(\ell, h, \iota)$ où
  $\ell, h \in \mathbb{R} \cup \{-\infty, +\infty\}$ et $\iota$ est un
  booléen (\code{is_int}). Sa concrétisation est
  \[
    \gam(\ell, h, \iota) \;=\; \{ x \in \mathbb{R} \mid \ell \le x \le h \}
    \;\cap\; \begin{cases} \Int & \text{si } \iota \\ \mathbb{R} & \text{sinon.}\end{cases}
  \]
  Tout triplet avec $\ell > h$ dénote $\emptyset$ : c'est $\Bot$. Le top est
  $[-\infty, +\infty]$ avec $\iota$ faux. La valeur JavaScript \code{NaN}
  n'appartient à la concrétisation d'\emph{aucun} intervalle ; les bornes
  infinies signifient \og pas de borne\fg{}.
\end{definition}

La dernière phrase est une décision, pas une évidence : elle oblige toute
opération qui peut produire \code{NaN} à sortir du domaine des intervalles
(\cref{sec:treillis-simples-arith}).

\subsection{Ce que fait le code}

\rustsnippet[label={lst:treillis-simples-interval}]{src/domains/impls/interval.rs}{5}{27}{la structure \code{Interval}}

\begin{itemize}
  \item \code{lo}, \code{hi} : les bornes, fermées, en \code{f64}.
    \code{f64::INFINITY} et \code{f64::NEG_INFINITY} codent l'absence de
    borne.
  \item \code{is_int} : \og toute valeur concrète est entière\fg{}. Vrai veut
    dire \emph{prouvé} entier ; faux veut dire \og peut contenir des
    non-entiers\fg{}. Le bit n'ajoute aucune valeur à la concrétisation, il
    en retire (les non-entiers) : c'est une annotation de précision.
  \item \code{PartialEq} est écrit à la main (l.~23--27) pour
    \emph{ignorer} \code{is_int} : deux intervalles de mêmes bornes sont
    égaux. Le point fixe converge ainsi quelle que soit l'évolution du bit.
\end{itemize}

\begin{invariant}{Sens du bit \code{is_int}}{treillis-simples-is-int}
  Si \code{is_int} est vrai, toute valeur concrète représentée est un entier.
  Toute opération qui produit un intervalle ne met \code{is_int} à vrai que si
  elle peut le prouver. Réciproquement, un \code{is_int} faux n'est jamais
  faux \og à tort\fg{} : il n'enlève rien.
\end{invariant}

Les constructeurs élémentaires :

\rustsnippet[label={lst:treillis-simples-interval-ctors}]{src/domains/impls/interval.rs}{29}{67}{constructeurs et tests élémentaires}

\begin{itemize}
  \item \code{point(v)} est l'abstraction d'une valeur unique. Son bit vaut
    vrai si $v$ est fini et sans partie fractionnaire : \code{point(3.0)} est
    entier, \code{point(1.5)} ne l'est pas, \code{point(f64::INFINITY)} non
    plus.
  \item \code{bottom()} est le $\Bot$ \emph{canonique} $[+\infty, -\infty]$,
    avec \code{is_int} vrai par vacuité (l'ensemble vide ne contient aucun
    non-entier).
  \item \code{is_bottom} teste $\ell > h$, et non l'égalité au $\Bot$
    canonique : c'est le seul test fiable (voir le piège plus bas).
\end{itemize}

\subsection{Ordre, join, meet}

\rustsnippet[label={lst:treillis-simples-interval-ord}]{src/domains/impls/interval.rs}{312}{328}{l'ordre des intervalles : l'inclusion des bornes}

L'ordre est l'inclusion : $[a,b] \lleq [c,d]$ ssi $c \le a$ et $b \le d$. Le
join est l'enveloppe convexe (\code{hull}, présentée au
\cref{chap:interpretation-abstraite},
\cref{lst:interpretation-abstraite-hull}) : les bornes extrêmes, entier si
les deux opérandes le sont, $\Bot$ neutre. L'implémentation du trait relie les
deux et définit le meet :

\rustsnippet[label={lst:treillis-simples-interval-impl}]{src/domains/impls/interval.rs}{332}{361}{\code{impl AbstractDomain for Interval} (début)}

Le meet est l'intersection des bornes. Son bit est un \emph{ou} : si l'un des
opérandes ne contient que des entiers, leur intersection non plus ; il n'y a
pas de test $\Bot$, car un résultat $\ell > h$ est $\Bot$ par construction.

\begin{proposition}{Treillis des intervalles}{treillis-simples-interval-treillis}
  Sur les intervalles non vides, \code{hull} est la plus petite borne
  supérieure et \code{meet} la plus grande borne inférieure pour l'ordre
  d'inclusion des bornes ; les deux respectent la concrétisation
  (conditions 3 et 6 de la \cref{def:treillis-simples-contrat}).
\end{proposition}
\begin{proof}
  Un intervalle $[c,d]$ majore $[a_1,b_1]$ et $[a_2,b_2]$ ssi
  $c \le \min(a_1,a_2)$ et $d \ge \max(b_1,b_2)$ ; le plus petit est
  $[\min(a_1,a_2), \max(b_1,b_2)]$, c'est \code{hull}. Si $x \in
  \gam(a_1,b_1,\iota_1)$, alors $\min(a_1,a_2) \le x \le \max(b_1,b_2)$ ; et si
  $\iota_1 \wedge \iota_2$, $x$ est entier. Le raisonnement est dual pour le
  meet : $x$ dans les deux concrétisations est entre les deux bornes basses et
  les deux bornes hautes, et entier dès que l'un des $\iota$ l'impose.
\end{proof}

\Cref{fig:treillis-simples-hasse-intervalles} montre un fragment de ce
treillis. Il est de hauteur infinie : la chaîne $[0,0] \sqsubset [0,1]
\sqsubset [0,2] \sqsubset \cdots$ ne s'arrête jamais, ce qui imposera le
widening (\cref{sec:treillis-simples-widening}).

\begin{figure}[htbp]\centering
\begin{tikzpicture}[treillis,xscale=1.2,yscale=1.1]
  \node (bot) at (0,0) {$\Bot$ \scriptsize($\ell > h$)};
  \node (p0) at (-2,1.2) {$[0,0]$};
  \node (p1) at (0,1.2) {$[1,1]$};
  \node (p2) at (2,1.2) {$[2,2]$};
  \node (pd) at (3.4,1.2) {$\cdots$};
  \node (pm) at (-3.4,1.2) {$\cdots$};
  \node (i01) at (-1,2.4) {$[0,1]$};
  \node (i12) at (1,2.4) {$[1,2]$};
  \node (i02) at (0,3.6) {$[0,2]$};
  \node (i03) at (-1.4,4.8) {$[0,3]$};
  \node (im12) at (1.4,4.8) {$[-1,2]$};
  \node (i0inf) at (-2.2,6.0) {$[0,+\infty]$};
  \node (top) at (0,7.2) {$[-\infty,+\infty] = \Top$};
  \foreach \n in {p0,p1,p2} { \draw (bot) -- (\n); }
  \draw[dashed,ra-gray] (bot) -- (pd) (bot) -- (pm);
  \draw (p0) -- (i01) -- (p1) -- (i12) -- (p2);
  \draw (i01) -- (i02) -- (i12);
  \draw (i02) -- (i03) (i02) -- (im12);
  \draw[dashed,ra-gray] (i03) -- (i0inf) (i0inf) -- (top) (im12) -- (top);
  \draw[flechemay] (p0.west) to[bend left=40] node[etiquette,left] {chaîne infinie} (i0inf.west);
\end{tikzpicture}
\caption{Un fragment du treillis des intervalles. Les points sont les
atomes ; le join de $[0,0]$ et $[2,2]$ est $[0,2]$, qui contient aussi $1$
(convexité). Toute chaîne $[0,0] \sqsubset [0,1] \sqsubset \cdots$ monte sans
fin vers $[0,+\infty]$ : hauteur infinie.}
\label{fig:treillis-simples-hasse-intervalles}
\end{figure}

\begin{exemple}{Le join perd les trous}{treillis-simples-hull}
  Un handler écrit \code{setX(0)}, un autre \code{setX(10)}. Le slot vaut
  $[0,0] \join [10,10] = [0,10]$, qui contient aussi $1, \dots, 9$. Le join est
  pourtant la plus petite borne supérieure \emph{dans ce treillis} : aucun
  intervalle plus petit ne contient $0$ et $10$. La perte vient du choix du
  domaine (convexe), pas de l'opération.
\end{exemple}

\begin{piege}
  $\Bot$ a plusieurs représentations. Le raffinement
  $[10,20]$ sous \code{x < 5} donne $[10, 4]$ (\cref{sec:treillis-simples-gardes}),
  un $\Bot$ non canonique. Les combinateurs (\code{hull}, \code{widen},
  \code{add}\dots) testent \code{is_bottom()} d'abord et le traitent
  correctement. Mais \code{PartialEq} et \code{PartialOrd} comparent les
  bornes brutes : \code{[10,4] == Interval::bottom()} est faux, et
  \code{partial_cmp([10,4], [0,3])} rend \code{None} (aucune des deux
  inclusions de bornes n'est vraie, \cref{lst:treillis-simples-interval-ord}),
  alors qu'un $\Bot$ devrait être \code{Less} que tout. Le
  \cref{chap:interpretation-abstraite} a relevé la conséquence pratique (une
  worklist qui revisite un bloc mort). Tout code nouveau qui compare des
  intervalles doit tester \code{is_bottom()} avant l'ordre.

  L'effet sur le point fixe est précis. \code{StateStore::leq}
  (\src{src/domains/stores/state_store.rs}{70}{77}) délègue à
  \code{leq_pointwise} (\src{src/domains/stores/mod.rs}{23}{25}), qui compare
  par \code{partial_cmp} brut ; \code{changed_labels}
  (\src{src/domains/stores/state_store.rs}{87}{95}) compare de même par
  \code{!=} brut. Si un slot contient un $\Bot$ non canonique dont la forme
  diffère d'un tour à l'autre (deux narrowings à des littéraux différents,
  par exemple), le test de convergence
  (\code{new_state.leq(&state)}, \src{src/engine/fixpoint.rs}{495}{495})
  peut rendre \code{false} alors que les deux valeurs dénotent
  $\emptyset$ : la boucle externe continue un tour de plus qu'il ne
  faudrait, et \code{changed_labels} peut marquer le slot \og changé\fg{} à
  tort, avançant son entrée dans \code{widen_trace}. Ce n'est ni un défaut de
  soundness (un tour de plus ne perd aucune valeur) ni un risque de
  non-terminaison : le garde-fou des 100 itérations
  (\src{src/engine/fixpoint.rs}{499}{512}) force de toute façon un
  \code{widen} final. L'impact observable est une perte de précision
  possible (un widening avancé d'un tour) sur un slot qui n'était de toute
  façon plus atteignable.
\end{piege}

\begin{piege}
  Parce que \code{is_int} est hors de l'égalité, un tour de point fixe qui
  n'ajouterait qu'un non-entier \emph{à l'intérieur} des bornes serait jugé
  convergé, et l'état retenu garderait un \code{is_int} vrai devenu faux.
  Le dossier de notes a tenté de reproduire ce scénario (un compteur par pas
  de $0{,}5$) sans y parvenir : le bit y est bien faux à la convergence. Le
  risque reste théorique, mais c'est la seule brèche de
  l'\cref{inv:treillis-simples-is-int}. Au commit \snapshotCommit{}, ni une
  reproduction ni une preuve d'impossibilité n'existent : le scénario exige
  que les bornes se stabilisent \emph{avant} qu'un opérande non entier
  n'entre en jeu à bornes inchangées, ce qu'aucun des exemples essayés
  (croissance directe, modulo) ne produit, faute d'avoir su construire un
  programme où l'un se produit sans l'autre.
\end{piege}

% ===========================================================================
\section{Arithmétique d'intervalles}
\label{sec:treillis-simples-arith}

\subsection{Addition, soustraction, multiplication, négation}

Chaque écriture d'état d'un compteur passe par une opération arithmétique :
\code{count + 1}, \code{total - price}, \code{n * 2}. Le domaine doit
calculer, à partir des intervalles des opérandes, un intervalle qui contient
tous les résultats possibles.

\rustsnippet[label={lst:treillis-simples-add-mul}]{src/domains/impls/interval.rs}{148}{187}{addition, soustraction, multiplication}

\rustsnippet[label={lst:treillis-simples-neg}]{src/domains/impls/interval.rs}{189}{198}{négation}

Ce sont les formules classiques de l'arithmétique d'intervalles :
\begin{align*}
  [a,b] + [c,d] &= [a+c,\; b+d], &
  [a,b] - [c,d] &= [a-d,\; b-c], \\
  [a,b] \times [c,d] &= [\min P,\; \max P],\ P = \{ac, ad, bc, bd\}, &
  -[a,b] &= [-b,\; -a].
\end{align*}
Un opérande $\Bot$ rend $\Bot$ : un chemin mort reste mort. Le bit entier est
un \emph{et} : la somme, la différence ou le produit de deux entiers est
entier, mais un seul non-entier suffit à perdre la garantie (test
\code{integrality_lost_through_float_arithmetic},
\src{src/domains/impls/interval.rs}{608}{613}).

\begin{theoreme}{Soundness de l'arithmétique d'intervalles}{treillis-simples-arith}
  Pour $x \in \gam(I)$ et $y \in \gam(J)$ finis, $x + y \in \gam(I + J)$,
  $x - y \in \gam(I - J)$, $-x \in \gam(-I)$, et, si aucun produit de coins
  n'est indéterminé ($0 \times \pm\infty$), $x y \in \gam(I \times J)$.
\end{theoreme}
\begin{proof}
  Addition : $a \le x \le b$ et $c \le y \le d$ donnent $a + c \le x + y \le
  b + d$. Soustraction : $-d \le -y \le -c$, puis addition. Négation :
  multiplier par $-1$ renverse les inégalités. Multiplication : pour $y$
  fixé, $x \mapsto xy$ est affine, donc ses extremums sur $[a,b]$ sont en
  $a$ ou $b$ ; puis $y \mapsto ay$ et $y \mapsto by$ sont affines, extremums
  en $c$ ou $d$ : les extremums de $xy$ sur le pavé sont parmi les quatre
  coins. Le bit : somme, différence, produit et opposé d'entiers sont
  entiers.
\end{proof}

Le calcul se fait en \code{f64}, avec l'arrondi au plus proche. Comme
JavaScript calcule les valeurs concrètes avec la même arithmétique et que
l'arrondi est monotone ($u \le v \Rightarrow \mathrm{rd}(u) \le
\mathrm{rd}(v)$), les bornes calculées encadrent encore les résultats
concrets. Cet argument n'est écrit nulle part dans le code ; il est
nécessaire, et il suffit.

La restriction \og aucun produit indéterminé\fg{} du théorème n'est pas une
précaution de style : le code la viole dans un cas, étudié à la
\cref{sec:treillis-simples-limites}.

\subsection{\code{NaN}, $\pm\infty$ et la division}

Puisque \code{NaN} n'est dans aucun intervalle, une opération qui peut le
produire doit rendre autre chose qu'un intervalle. La division en est
l'exemple type : \code{x / 0} vaut $\pm$\code{Infinity}, et \code{0 / 0} vaut
\code{NaN}. Le code ne cherche pas à distinguer les cas : \code{BinOp::Div}
rend toujours $\Top$ au niveau \code{StateValue}
(\cref{lst:treillis-simples-eval-binop}). \code{docs/limitations.md} le
documente.

La place de $\pm\infty$ est plus ambiguë. Une borne infinie signifie
\og non borné\fg{}, mais un littéral \code{1e999} vaut concrètement
\code{Infinity}, et \code{point(f64::INFINITY)} produit l'intervalle
$[+\infty, +\infty]$. Le code traite de fait les infinis comme des
\emph{absences de borne}, pas comme des valeurs : les formules ci-dessus
donneraient $+\infty - (+\infty) = \code{NaN}$ comme borne si les deux
opérandes étaient le point $+\infty$. Aucun programme React raisonnable ne
fait d'arithmétique sur \code{Infinity}, et l'écart n'a pas été observé sur
un diagnostic ; il est signalé à la \cref{sec:treillis-simples-limites}.

\subsection{Reste et puissance (commit \code{548f922})}

Jusqu'au commit \code{548f922} (issue \issue{73}), \texttt{\%} et \code{**}
n'avaient pas de variante propre dans l'IR et tombaient dans l'opérateur
opaque. Ils sont maintenant modélisés :

\rustsnippet[label={lst:treillis-simples-rem-pow}]{src/domains/impls/interval.rs}{200}{253}{reste et puissance}

Les deux fonctions rendent une \code{Option} : \code{None} signifie
\og le résultat peut être \code{NaN}\fg{}, et l'appelant le transforme en
$\Top$.

\paragraph{Le reste.} En JavaScript, \texttt{x \% m} a le signe du dividende
\code{x}, et sa valeur absolue est strictement inférieure à $|m|$ et
inférieure ou égale à $|x|$. Le code en tire :
\begin{itemize}
  \item si le diviseur peut être nul ($c \le 0 \le d$), \code{None} ;
  \item sinon, $\mu = \max(|c|, |d|)$, la plus grande valeur absolue
    possible du diviseur, et la borne $\beta = \mu - 1$ sur des entiers
    (strictement sous $\mu$), $\beta = \mu$ sinon (borne fermée, donc
    sur-approximation d'une borne ouverte) ;
  \item la borne basse est $0$ si le dividende est positif, sinon
    $\max(a, -\beta)$ ; la borne haute est $0$ si le dividende est négatif,
    sinon $\min(b, \beta)$.
\end{itemize}

\begin{proposition}{Soundness du reste}{treillis-simples-rem}
  Si $0 \notin [c,d]$, pour tous $x \in \gam([a,b])$, $m \in \gam([c,d])$
  finis, $x \mathbin{\%} m \in \gam(\code{rem}([a,b],[c,d]))$.
\end{proposition}
\begin{proof}
  Soit $r = x \mathbin{\%} m$. Si $x \ge 0$, $0 \le r \le x \le b$ et
  $r < |m| \le \mu$ ; sur des entiers $r \le \mu - 1$. Donc
  $r \le \min(b, \beta)$ et $r \ge 0 \ge$ la borne basse (qui vaut $0$ ou
  $\max(a,-\beta) \le 0$). Le cas $x \le 0$ est symétrique. Si $a \ge 0$, tout
  $x$ est positif et $r \ge 0$ : la borne basse $0$ est correcte ; si
  $b \le 0$, la borne haute $0$ l'est de même. Le bit : le reste de deux
  entiers est entier.
\end{proof}

\begin{exemple}{Un compteur modulo}{treillis-simples-mod}
  \texttt{setI((i + 1) \% 3)} à partir de \code{i = 0}. Au premier tour,
  $[0,0] + [1,1] = [1,1]$ et $\code{rem}([1,1],[3,3]) = [0, \min(1, 2)] =
  [0,1]$ ; le slot vaut $[0,0] \join [0,1] = [0,1]$. Au deuxième,
  $[1,2] \mathbin{\%} [3,3] = [0,2]$ ; au troisième, $[1,3] \mathbin{\%} [3,3]
  = [0,2]$ : point fixe, sans widening. Avant le commit \code{548f922}, le
  résultat était $\Top$.
\end{exemple}

\paragraph{La puissance.} Sur le quadrant $x \ge 0, y \ge 0$, la fonction
$x^y$ est croissante en $x$ pour $y$ fixé, et monotone en $y$ pour $x$ fixé
(croissante si $x > 1$, constante si $x = 1$, décroissante si $x < 1$). Ses
extremums sur un pavé sont donc aux coins, comme pour le produit. Hors du
quadrant, un exposant fractionnaire d'une base négative donne \code{NaN}, un
exposant négatif de zéro donne \code{Infinity} : \code{None}. On note que
$0^0 = 1$ en JavaScript comme pour \code{f64::powf}.

\begin{decision}[commit 548f922 --- une variante par opérateur]
  Le commit qui modélise \texttt{\%} et \code{**} applique une règle prise avec
  \issue{74} et \issue{75} : chaque opérateur du langage a sa variante de
  \code{BinOp}, jamais un opérateur source porté par \code{Unknown}. Avec
  \code{Mod}, \code{Pow}, \code{In} et \code{InstanceOf}, \code{BinOp::Unknown}
  n'a plus de producteur et disparaît ; \code{lower_binop} devient exhaustif,
  si bien qu'un opérateur nouveau du langage ne peut plus tomber
  silencieusement dans l'opaque. L'alternative (garder un joker \code{_}
  vers $\Top$) était sound mais masquait les oublis. Le corpus, à 1494
  diagnostics avant le commit, n'a pas bougé (\file{docs/corpus-baseline.json}
  n'est pas modifié par ce commit).
\end{decision}

\subsection{La table des opérateurs binaires}

Les intervalles ne sont jamais utilisés seuls : l'évaluateur manipule des
\code{StateValue}, et il décide, opérateur par opérateur, s'il peut se ramener
aux intervalles.

\rustsnippet[label={lst:treillis-simples-eval-binop}]{src/domains/transfer/state_value.rs}{752}{811}{la table des opérateurs binaires}

Le passage aux intervalles se fait par \code{as_arith}
(\src{src/domains/transfer/state_value.rs}{708}{732}), qui rend un intervalle
si les seules sortes possibles de la valeur sont \og nombre\fg{} et
\og \code{null}\fg{}, puisque \code{ToNumber(null)} vaut $0$ ; le
\cref{chap:statevalue-stabilite} en détaille la logique. Il rend \code{None}
dès qu'un \code{undefined}, un booléen, une chaîne, une référence, un setter
ou la sorte résiduelle \code{other} est possible. La lecture de la table,
ligne par ligne :
\begin{itemize}
  \item \code{+} : numérique si les deux opérandes le sont ; concaténation
    exacte si les deux sont des chaînes pures ; \og une chaîne\fg{} si ce sont
    des chaînes mal connues ; $\Top$ sinon. \code{1 + "a"} donne donc
    $\Top$ et non une chaîne : sûr, imprécis.
  \item \code{-} et \code{*} : numériques ou $\Top$ ; \code{/} : toujours
    $\Top$ ; \texttt{\%} et \code{**} : l'intervalle, ou $\Top$ si \code{None}.
  \item \code{&&} et \code{||} : $\Top$. Leur valeur est rarement utilisée,
    car le lowering les transforme en diamant de CFG
    (\cref{chap:lowering-cfg}).
  \item comparaisons, \code{in}, \code{instanceof} : un booléen inconnu
    (\cref{sec:treillis-simples-plats}).
  \item bit à bit : \code{eval_bitwise}, section suivante.
\end{itemize}

% ===========================================================================
\section{Les opérateurs bit à bit}
\label{sec:treillis-simples-bitwise}

Le code React utilise les opérateurs bit à bit pour des drapeaux
(\code{flags & MASK}), des hachages (\texttt{hash >{}>{}> 0}) ou des index
(\texttt{i <{}< 1}). Jusqu'à \issue{74}, ils valaient $\Top$ ; ils sont
maintenant modélisés à partir d'une observation simple.

\begin{remarque}
  JavaScript convertit les deux opérandes d'un opérateur bit à bit par
  \code{ToInt32} (\code{ToUint32} pour l'opérande gauche de \texttt{>{}>{}>}) :
  troncature vers zéro, puis réduction modulo $2^{32}$ dans
  $[-2^{31}, 2^{31}-1]$. Le résultat est donc \emph{toujours} un entier de
  cette plage, quels que soient les opérandes (chaînes, objets,
  \code{undefined}\dots), et un entier de $[0, 2^{32}-1]$ pour \texttt{>{}>{}>}.
\end{remarque}

\rustsnippet[label={lst:treillis-simples-int32}]{src/domains/transfer/state_value.rs}{813}{844}{plages int32/uint32 et utilitaires}

\rustsnippet[label={lst:treillis-simples-bitwise-a}]{src/domains/transfer/state_value.rs}{846}{886}{\code{eval_bitwise} : plancher, \code{&}, \code{|} et \code{^}}

\rustsnippet[label={lst:treillis-simples-bitwise-b}]{src/domains/transfer/state_value.rs}{887}{939}{\code{eval_bitwise} : décalages, et \code{all_ones_above}}

Le principe : un résultat \emph{plancher} (\code{int32_range} ou
\code{uint32_range}, l.~925--928), qui ne demande aucune hypothèse sur les
opérandes, et un résultat \emph{raffiné} quand un opérande est constant ou
borné. Rendre la plage int32 plutôt que $\Top$ est déjà un gain : les
composantes chaîne, booléen et référence restent à $\Bot$, et une garde sur
le résultat reste raffinable. Cas par cas :
\begin{itemize}
  \item \code{x & m}, avec $m$ une constante de $[0, 2^{31}-1]$ (à gauche ou
    à droite) : le résultat n'a que des bits de $m$, donc est dans $[0, m]$.
  \item \code{x | y}, \code{x ^ y}, les deux opérandes positifs dans int32 :
    le résultat ne peut pas avoir de bit au-dessus du plus haut bit possible
    des opérandes, donc il est dans $[0, 2^n - 1]$ où $2^n - 1$ est le plus
    petit nombre de la forme \og que des 1\fg{} au moins égal au max
    (\code{all_ones_above}).
  \item \texttt{x <{}< k}, $k$ constant : multiplication des bornes par $2^k$,
    gardée seulement si le résultat ne déborde pas d'int32 (au-delà, la
    réduction modulo $2^{32}$ rend l'opération non monotone).
  \item \texttt{x >{}> k} : division entière par $2^k$ arrondie vers le bas,
    monotone.
  \item \texttt{x >{}>{}> k} : le résultat tient dans les $32 - k$ bits bas, quel
    que soit $x$.
  \item Un opérande $\Bot$ rend $\Bot$ (l.~857--859).
\end{itemize}

La négation bit à bit \code{~x}, qui vaut
$-(\code{ToInt32}(x) + 1)$, est traitée à part dans les opérateurs unaires,
avec soin :

\rustsnippet[label={lst:treillis-simples-bitnot}]{src/domains/transfer/state_value.rs}{978}{993}{la négation bit à bit}

La fonction est décroissante, les bornes s'échangent, et les arrondis
\code{ceil} sur la borne haute et \code{floor} sur la borne basse encadrent
la troncature de \code{ToInt32} sur des non-entiers : pour $x = 1{,}5$,
$\code{ToInt32}(x) = 1$ et $\mathord{\sim} x = -2$, dans
$[-\lceil 1{,}5\rceil - 1, -\lfloor 1{,}5 \rfloor - 1] = [-3, -2]$.

\begin{exemple}{Calculs bit à bit}{treillis-simples-bitwise}
  $[0,100] \mathbin{\&} 15 = [0,15]$. $[0,5] \mathbin{|} [0,9]$ : le max est
  $9$, $\lfloor \log_2 9 \rfloor + 1 = 4$ bits, résultat $[0, 15]$.
  $[3,3] \ll 2 = [12,12]$. $[0,1000] \gg 2 = [0, 250]$.
  $x \ggg 28 \in [0, 15]$ quel que soit $x$.
  $\mathord{\sim}[-3,4] = [-5, 2]$.
\end{exemple}

\subsection{Un faux négatif : les décalages sur des non-entiers}

Les bras \texttt{<{}<} et \texttt{>{}>} supposent implicitement que les bornes de
l'opérande sont des entiers : ils multiplient ou divisent \code{a.lo} et
\code{a.hi} tels quels. Or \code{ToInt32} \emph{tronque} d'abord. Pour
$x = 1{,}5$, JavaScript calcule $1{,}5 \ll 1 = \code{ToInt32}(1{,}5) \ll 1 =
1 \ll 1 = 2$, alors que le code calcule $[1{,}5 \times 2, 1{,}5 \times 2] =
[3,3]$, qui ne contient pas $2$. De même pour $x = -1{,}5$ :
$-1{,}5 \gg 0 = \code{ToInt32}(-1{,}5) = -1$, alors que le code calcule
$[\lfloor -1{,}5 \rfloor, \lfloor -1{,}5 \rfloor] = [-2, -2]$. Ce sont des
sous-approximations, la direction interdite.

Le défaut est observable de bout en bout :

\begin{lstlisting}[language=TypeScript,caption={Décalages sur un état non entier (\file{/tmp/ch09/shift2.tsx}, élagué de sa ligne d'import)},label={lst:treillis-simples-shift}]
export function ShlFloat() {
  const [x] = useState(1.5);
  const [count, setCount] = useState(0);
  useEffect(() => {
    const y = x << 1;
    setCount(count + 3 - y);
  }, [count]);
  return <div>{count}</div>;
}

export function ShrFloat() {
  const [x] = useState(-1.5);
  const [count, setCount] = useState(0);
  useEffect(() => {
    const y = x >> 0;
    setCount(count + y + 2);
  }, [count]);
  return <div>{count}</div>;
}

export function ShlInt() {
  const [x] = useState(1);
  const [count, setCount] = useState(0);
  useEffect(() => {
    const y = x << 1;
    setCount(count + 3 - y);
  }, [count]);
  return <div>{count}</div>;
}
\end{lstlisting}

Concrètement, les trois composants bouclent : dans \code{ShlFloat},
\code{y} vaut $2$ et l'effet écrit \code{count + 1} ; dans \code{ShrFloat},
\code{y} vaut $-1$ et l'effet écrit encore \code{count + 1} ; \code{ShlInt}
est le témoin entier. L'analyseur :

\begin{sortie}
$ reactant check /tmp/ch09/shift2.tsx --show-clean
  ShlFloat  (3 hooks)  /tmp/ch09/shift2.tsx  ✓
  ShlInt  (3 hooks)  /tmp/ch09/shift2.tsx
    warn   infinite-loop  [hook:1]  (line 26:2)  this effect keeps pushing state `count` (its deps do not provably gate it, so the effect can re-run every render) to new values on every run. Potential infinite render loop
       (2 trace step(s), rerun with --trace)
  ShrFloat  (3 hooks)  /tmp/ch09/shift2.tsx  ✓

⚠  1 warning(s) across 1 file(s).
\end{sortie}

Pour l'analyseur, \code{y} vaut $[3,3]$ dans \code{ShlFloat} et $[-2,-2]$
dans \code{ShrFloat} : l'effet écrit \code{count + 0}, la valeur ne change
pas, le point fixe converge sans widening, et les deux boucles infinies sont
déclarées sans défaut. C'est un faux négatif ; aucune issue ne le suit au
commit étudié. La correction naturelle, centrale et d'une ligne par bras,
consiste à tronquer les bornes avant de décaler (\code{trunc} est monotone,
et coïncide avec \code{ToInt32} sur la plage int32), ou à exiger
\code{is_int} ; c'est l'\cref{exo:treillis-simples-shift}.

% ===========================================================================
\section{Widening : du saut à l'infini aux seuils}
\label{sec:treillis-simples-widening}

\subsection{La chronologie d'un compteur}

Le treillis des intervalles est de hauteur infinie : sans précaution, le
point fixe de \code{Unbounded} ne termine pas. Suivons la boucle externe du
moteur (\cref{chap:point-fixe-composant}) sur les deux compteurs du
\cref{lst:treillis-simples-counter}. Le slot est amorcé à $[0,0]$ par
\code{useState(0)}. À chaque tour, l'effet lit la valeur courante, écrit
\code{count + 1}, et l'écriture est ajoutée au slot par join (mise à jour
faible). Tant que le compteur de tours \code{iteration} reste sous
\code{widen_threshold} (3 par défaut, \src{src/engine/fixpoint.rs}{51}{59}),
le nouvel état remplace l'ancien ; à partir de 3, il est \emph{élargi}
(\src{src/engine/fixpoint.rs}{514}{529}, reproduit au
\cref{lst:interpretation-abstraite-outer-widen}) :

\begin{center}
\code{state = state.widen_to(&new_state, &thresholds)}
\end{center}

\Cref{fig:treillis-simples-chronologie} dessine les deux histoires.

\begin{figure}[htbp]\centering
\begin{tikzpicture}[x=26mm,y=11mm]
  \foreach \i/\t in {0/{amorçage},1/{tour 1},2/{tour 2},3/{tour 3 (widening)},4/{tour 4}} {
    \node[etiquette] at (\i,0.9) {\t};
  }
  \draw[fleche] (-0.4,0.45) -- (4.5,0.45);
  % Unbounded
  \node[etiquette,anchor=east] at (-0.35,-0.3) {\code{Unbounded}};
  \node[bloc] (u0) at (0,-0.3) {$[0,0]$};
  \node[bloc] (u1) at (1,-0.3) {$[0,1]$};
  \node[bloc] (u2) at (2,-0.3) {$[0,2]$};
  \node[bloc,draw=ra-amber] (u3) at (3,-0.3) {$[0,+\infty]$};
  \node[bloc] (u4) at (4,-0.3) {$[0,+\infty]$ \checkmark};
  \draw[fleche] (u0) -- (u1); \draw[fleche] (u1) -- (u2); \draw[fleche] (u2) -- (u3); \draw[fleche] (u3) -- (u4);
  \node[etiquette] at (3,-1.05) {$[0,2] \widen_{\{0,1\}} [0,3]$};
  % Guarded
  \node[etiquette,anchor=east] at (-0.35,-2.3) {\code{Guarded}};
  \node[bloc] (g0) at (0,-2.3) {$[0,0]$};
  \node[bloc] (g1) at (1,-2.3) {$[0,1]$};
  \node[bloc] (g2) at (2,-2.3) {$[0,2]$};
  \node[bloc,draw=ra-amber] (g3) at (3,-2.3) {$[0,10]$};
  \node[bloc] (g4) at (4,-2.3) {$[0,10]$ \checkmark};
  \draw[fleche] (g0) -- (g1); \draw[fleche] (g1) -- (g2); \draw[fleche] (g2) -- (g3); \draw[fleche] (g3) -- (g4);
  \node[etiquette] at (3,-3.05) {$[0,2] \widen_{\{0,1,10\}} [0,3]$};
  \node[etiquette,align=center] at (4,-3.85) {lecture $[0,10]$, garde $\Rightarrow [0,9]$,\\écriture $[1,10]$ : stable};
\end{tikzpicture}
\caption{Chronologie du slot \code{count} dans la boucle externe. Aux tours
1 et 2, le nouvel état remplace l'ancien ; au tour 3 (\code{iteration = 3}),
il est élargi à seuils : la borne haute saute au plus petit seuil
$\ge 3$, qui est $+\infty$ pour \code{Unbounded} ($T = \{0,1\}$) et $10$ pour
\code{Guarded} ($T = \{0,1,10\}$). Au tour 4, le nouvel état est inclus dans
l'ancien : post-point fixe (\checkmark).}
\label{fig:treillis-simples-chronologie}
\end{figure}

Détaillons \code{Guarded}. Tour 1 : l'effet lit $[0,0]$, la garde
\code{count < 10} le laisse intact, l'écriture vaut $[1,1]$, le slot
$[0,0] \join [1,1] = [0,1]$. Tour 2 : lecture $[0,1]$, écriture $[1,2]$, slot
$[0,2]$. Tour 3 : l'état candidat est $[0,3]$ ; comme \code{iteration}
atteint 3, on calcule $[0,2] \widen_T [0,3]$ avec $T = \{0, 1, 10\}$ (les
littéraux du composant) : la borne haute a grandi, elle prend le plus petit
seuil $\ge 3$, soit $10$. Tour 4 : la lecture vaut $[0,10]$, la garde la
raffine en $[0,9]$ (\cref{sec:treillis-simples-gardes}), l'écriture vaut
$[1,10]$, le candidat $[0,10]$ est inclus dans l'état : c'est fini. Les
valeurs finales $[0,+\infty]$ et $[0,10]$ sont épinglées par les tests
\code{unbounded_counter_is_flagged_and_grows_to_infinity} et
\code{guarded_counter_converges_bounded}
(\src{tests/widening_e2e.rs}{79}{122}). La sortie \texttt{-{}-info} en
garde la trace :

\begin{sortie}
$ reactant check /tmp/ch09/counter.tsx --info --trace
  Guarded  (2 hooks)  /tmp/ch09/counter.tsx
    info   widening-info  (line 13:2)  state `count` kept changing during analysis and was approximated to converge, so findings that depend on it may be imprecise
       → state `count` is written here [hook:1] (line 13:2)
       → the abstract value of state `count` kept growing and was widened at iteration 3
  Unbounded  (2 hooks)  /tmp/ch09/counter.tsx
    warn   infinite-loop  [hook:0]  (line 5:2)  this effect keeps pushing state `count` (its deps do not provably gate it, so the effect can re-run every render) to new values on every run. Potential infinite render loop
       → state `count` is written here [hook:1] (line 5:2)
       → the abstract value of state `count` kept growing and was widened at iteration 3
    info   widening-info  (line 5:2)  state `count` kept changing during analysis and was approximated to converge, so findings that depend on it may be imprecise
       → state `count` is written here [hook:1] (line 5:2)
       → the abstract value of state `count` kept growing and was widened at iteration 3

⚠  1 warning(s) across 1 file(s).
\end{sortie}

(Sortie élaguée des lignes \texttt{verified} : \texttt{--info} liste aussi,
pour chaque composant, les règles qui ne signalent rien.) Les deux slots ont
été élargis au tour 3 ; seul celui d'\code{Unbounded} est
non borné, et c'est ce que lit \regle{infinite-loop}.

\subsection{Les deux opérateurs}

Le widening simple (\src{src/domains/impls/interval.rs}{85}{106}) et le
widening à seuils (\src{src/domains/impls/interval.rs}{108}{146}) ont été
reproduits au \cref{chap:interpretation-abstraite}
(\cref{lst:interpretation-abstraite-widen,lst:interpretation-abstraite-widen-to}),
avec leur formule (\cref{def:widening}). Rappelons celle du second : une
borne qui n'a pas grandi est conservée ; une borne haute qui a grandi jusqu'à
$d$ devient $\min\{t \in T \cup \{+\infty\} \mid t \ge d\}$, une borne basse
qui a baissé jusqu'à $c$ devient $\max\{t \in T \cup \{-\infty\} \mid t \le
c\}$. Le code calcule ces extremums par \code{filter} puis \code{fold} à
partir de $\pm\infty$, sans exiger que $T$ soit trié.

\begin{theoreme}{\code{widen_to} est un widening}{treillis-simples-widen-to}
  Pour tout ensemble fini $T \subseteq \mathbb{R}$ :
  (i) $I \join J \lleq I \widen_T J$ pour tous $I, J$ ;
  (ii) pour toute suite $J_1, J_2, \dots$, la suite $I_0$,
  $I_{n+1} = I_n \widen_T J_{n+1}$ est ultimement stationnaire.
\end{theoreme}
\begin{proof}
  (i) Si $I$ ou $J$ est $\Bot$, le code rend l'autre, qui est le join. Sinon,
  la borne haute du résultat est $h_I$ si $h_J \le h_I$, et un $t \ge h_J$
  sinon (seuil ou $+\infty$) : dans les deux cas elle est
  $\ge \max(h_I, h_J)$. Symétriquement pour la borne basse. Le bit est
  $\iota_I \wedge \iota_J$, comme le join. (ii) La borne haute de $I_n$ ne
  décroît jamais : elle est conservée ou remplacée par un $t \ge h_{J} > h_{I_n}$.
  Après son premier saut, elle appartient à $T \cup \{+\infty\}$, ensemble
  fini, et chaque saut ultérieur la fait croître strictement dans cet
  ensemble : au plus $|T| + 1$ sauts. Idem pour la borne basse. Une suite
  dont les deux bornes cessent de changer est stationnaire.
\end{proof}

Le test \code{widen_to_is_sound_superset_of_hull}
(\src{src/domains/impls/interval.rs}{568}{576}) vérifie (i) sur un exemple.
Le point (ii) explique pourquoi l'ensemble des seuils doit être \emph{fini}
et \emph{fixé} avant le point fixe : des seuils qui dépendraient des valeurs
calculées pourraient ne jamais s'épuiser.

\begin{remarque}
  \code{widen_to} ne s'applique qu'à la composante \code{num} du produit.
  \code{StateValue::widen_to} le dit en deux lignes : seuils pour
  \code{num}, widening ordinaire (donc join, pour les composantes de hauteur
  finie) pour le reste.
\end{remarque}

\rustsnippet[label={lst:treillis-simples-sv-widen-to}]{src/domains/impls/state_value.rs}{670}{676}{le widening à seuils du produit}

\subsection{D'où viennent les seuils}

\code{collect_thresholds} (\src{src/engine/fixpoint.rs}{1185}{1206},
reproduit au \cref{lst:interpretation-abstraite-thresholds}) récolte une
fois par composant les littéraux numériques du rendu, des initialiseurs de
\code{useState}, des corps d'effets et de handlers. La descente dans les
expressions est la suivante :

\rustsnippet[label={lst:treillis-simples-collect-lits}]{src/engine/fixpoint.rs}{1212}{1224}{la récolte des littéraux, fermetures comprises}

Le cas \code{FnLit} est notable : la descente structurelle ordinaire ne
traverse pas les fermetures, mais un seuil écrit dans une fermeture
(\code{onClick={() => n < 5 && setN(n + 1)}}) doit pouvoir borner une
croissance gardée. Le résultat est trié et dédupliqué. Pour \code{Unbounded},
$T = \{0, 1\}$ (l'initialiseur et l'incrément) ; pour \code{Guarded},
$T = \{0, 1, 10\}$.

\begin{piege}
  Tout littéral numérique devient un seuil, même sans rapport avec la
  croissance étudiée. Un \code{padding: 12} dans le rendu d'un compteur non
  gardé ferait sauter la borne à $12$ au tour 3, puis, au tour suivant, à
  $+\infty$ (aucun seuil $\ge 13$). La précision peut changer, jamais la
  soundness (théorème ci-dessus), mais un seuil de plus peut retarder la
  convergence d'un tour, et changer le numéro d'itération qu'affiche
  \regle{widening-info}.
\end{piege}

\subsection{Le même widening dans les boucles locales}

La boucle externe n'est pas le seul endroit où une valeur peut croître sans
fin : une boucle \code{while} dans un effet en est un autre. \code{analyze_cfg}
applique le même opérateur, avec les mêmes seuils, sur les arcs arrière du
CFG (\src{src/engine/cfg_analyzer.rs}{97}{112}), à partir du troisième passage
par un même en-tête (\cref{chap:cfg-analyzer-dominance}). Sur le composant
\code{BoundedLocalLoop} de \file{tests/fixtures/widening.tsx}
(\code{let i = 0; while (i < 5) { i = i + 1; } setTotal(i);}), l'en-tête de
boucle voit $[0,0]$, $[0,1]$, $[0,2]$, puis $[0,2] \widen_{\{0,1,5\}} [0,3] =
[0,5]$ ; le corps, raffiné par \code{i < 5}, écrit $[1,5]$, et l'en-tête est
stable. La sortie de boucle, raffinée par la négation \code{i >= 5}, vaut
$[5,5]$, et le slot \code{total} converge vers $[0,5]$
(test \code{bounded_local_loop_is_precise},
\src{tests/widening_e2e.rs}{125}{137}).

\begin{decision}[ADR-14 --- seuils, signature inchangée, pas de narrowing]
  \adr{14} ajoute le widening à seuils \og à la ASTRÉE\fg{} sur un ensemble
  fini récolté une fois par composant, appliqué à la boucle externe et aux
  arcs arrière. Deux choix de détail sont ceux du code actuel. D'abord, la
  signature de \code{AbstractDomain::widen} n'est pas modifiée (la faire
  porter des seuils aurait fait fuir un paramètre inutile dans tous les
  domaines et les stores) : on ajoute \code{widen_to}, avec une
  implémentation par défaut qui ignore les seuils. Ensuite, la partie 2 de
  l'ADR (un opérateur \code{narrow} et une phase descendante bornée) est
  abandonnée par la révision du 28 juin 2026 : le raffinement sur garde
  pendant la montée et les seuils récupèrent déjà les bornes littérales, et
  une phase descendante ne pourrait pas faire redescendre le store d'état,
  où les écritures s'accumulent par join. L'ADR prévoyait aussi de ne
  récolter que les littéraux des gardes et des initialiseurs ; le code en
  récolte davantage, ce qui est sans danger. Les éléments annexes de l'ADR
  (un type \code{WidenOutcome}) n'ont pas été réalisés : le nom n'apparaît
  nulle part dans \file{src/}.
\end{decision}

\begin{piege}
  Le mot \og narrowing\fg{} a deux sens. En théorie de l'interprétation
  abstraite, c'est l'opérateur de la phase descendante, que \reactant{}
  n'implémente pas. Dans le code, les méthodes \code{narrow_*} sont des
  \emph{raffinements sur garde} : elles restreignent une variable sous une
  condition de branche, pendant la montée. C'est l'objet de la section
  suivante. \adr{14} insiste sur la distinction.
\end{piege}

% ===========================================================================
\section{Raffinement sur gardes}
\label{sec:treillis-simples-gardes}

\subsection{Le besoin}

Sans raffinement, la garde de \code{Guarded} ne sert à rien : l'effet lit
$[0,10]$, écrit $[1,11]$, et le widening fait sauter la borne à $+\infty$.
Il faut que, \emph{dans la branche} \code{if (count < 10)}, la variable
\code{count} soit restreinte aux valeurs qui satisfont la condition. C'est
le rôle des \code{narrow_*}. Leur contrat est la condition 5 de la
\cref{def:treillis-simples-contrat} : garder au moins toutes les valeurs
concrètes qui satisfont la condition.

\subsection{Ce que fait le code : les raffinements numériques}

\rustsnippet[label={lst:treillis-simples-narrow}]{src/domains/impls/interval.rs}{255}{309}{les raffinements numériques des intervalles}

Pour une comparaison stricte à un littéral $v$, la borne sûre sur les réels
est $v$ elle-même, fermée : $x < v \Rightarrow x \le v$ est vrai, et l'on ne
peut pas représenter une borne ouverte. Sur des entiers, en revanche,
$x < v \Rightarrow x \le \lceil v \rceil - 1$. Le code applique cette
seconde forme seulement si \code{is_int} est prouvé. C'est tout l'intérêt
du bit : sans lui, $[0,10]$ sous \code{count < 10} resterait $[0,10]$,
l'écriture vaudrait $[1,11]$, et le compteur gardé ne convergerait pas vers
une borne finie.

\begin{theoreme}{Soundness des raffinements numériques}{treillis-simples-narrow}
  Pour tout intervalle $I$ et tout littéral $v$ :
  $\gam(I) \cap \{x \mid x < v\} \subseteq \gam(\code{narrow_lt}(I, v))$,
  et de même pour \code{narrow_leq}, \code{narrow_gt}, \code{narrow_geq},
  \code{narrow_eq}, \code{narrow_neq} avec leurs conditions respectives.
\end{theoreme}
\begin{proof}
  \code{narrow_lt} : soit $x \in \gam(I)$ avec $x < v$. Si $I$ n'est pas
  entier, $x \le v$ et $x \le h$, donc $x \le \min(h, v)$. Si $I$ est
  entier, $x$ est un entier $< v$, donc $x \le \lceil v \rceil - 1$ (pour
  $v$ entier, $\lceil v \rceil - 1 = v - 1$ ; pour $v$ non entier, le plus
  grand entier $< v$ est $\lfloor v \rfloor = \lceil v \rceil - 1$). La
  borne basse et le bit ne changent pas, et $x$ reste entier. Les cas
  \code{narrow_leq} ($\lfloor v \rfloor$), \code{narrow_gt}
  ($\lfloor v \rfloor + 1$), \code{narrow_geq} ($\lceil v \rceil$) sont
  symétriques. \code{narrow_eq} : si $v \in [\ell, h]$, le seul $x$ possible
  est $v$, et \code{point(v)} le contient ; sinon aucun $x$ ne convient et
  $\Bot$ est exact. \code{narrow_neq} : rendre $I$ est toujours sûr ; rendre
  $\Bot$ quand $I$ est le point $v$ l'est aussi, car aucun $x \neq v$ n'est
  dans $\gam([v,v])$.
\end{proof}

\begin{exemple}{Arrondis entiers}{treillis-simples-arrondis}
  Sur un intervalle entier $[0,9]$ : sous \code{x < 4.5}, borne
  $\lceil 4{,}5 \rceil - 1 = 4$, résultat $[0,4]$ ; sous \code{x <= 4.5},
  $\lfloor 4{,}5 \rfloor = 4$ ; sous \code{x > 4.5}, $\lfloor 4{,}5\rfloor +
  1 = 5$, résultat $[5,9]$. Sur le point non entier $[1{,}7; 1{,}7]$ sous
  \code{x < 2}, le bit est faux, la borne est $2$, et $1{,}7$ est gardé. Un
  code qui aurait appliqué $x < 2 \Rightarrow x \le 1$ sans condition aurait
  rendu $\Bot$ et perdu une valeur réelle ; c'était le cas avant
  l'introduction du bit \code{is_int}, pour corriger précisément ce faux
  négatif (tests \code{narrow_lt_float_keeps_boundary_value} et
  \code{narrow_lt_integer_still_tightens},
  \src{src/domains/impls/interval.rs}{587}{606}).
\end{exemple}

\code{narrow_neq} est volontairement faible : un intervalle ne peut pas être
coupé en son milieu ($[0,5]$ privé de $3$ n'est pas convexe), si bien que
seul le cas d'un point égal à $v$ est exploité. Il en va de même pour les
bornes : $[0,5]$ privé de $0$ \emph{pourrait} devenir $[1,5]$ sur des
entiers, mais le code ne le fait pas (précision non implémentée, sans
conséquence sur la soundness).

\subsection{Qui appelle les raffinements}

Le moteur choisit l'opérateur à partir de la \emph{forme syntaxique} de la
condition, dans \code{narrow_env_for_branch}
(\cref{chap:cfg-analyzer-dominance}). Pour les comparaisons numériques :

\rustsnippet[label={lst:treillis-simples-narrow-env}]{src/engine/cfg_analyzer.rs}{212}{241}{choix du raffinement numérique selon la condition et la branche}

Le paramètre \code{taken} distingue la branche \og alors\fg{} (condition vraie)
de la branche \og sinon\fg{} (condition fausse, donc négation). La négation
de \code{x < v} est \code{x >= v}, d'où les paires croisées. Le
\cref{tab:treillis-simples-motifs} résume tous les motifs reconnus.

\begin{table}[htbp]\centering\small
\begin{tabular}{lll}
\hline
Condition & branche vraie & branche fausse \\
\hline
\code{x < v} / \code{x <= v} & \code{narrow_lt} / \code{narrow_leq} & \code{narrow_geq} / \code{narrow_gt} \\
\code{x > v} / \code{x >= v} & \code{narrow_gt} / \code{narrow_geq} & \code{narrow_leq} / \code{narrow_lt} \\
\code{x == v}, \code{x === v} & \code{narrow_eq} & \code{narrow_neq} \\
\code{x != v}, \code{x !== v} & \code{narrow_neq} & \code{narrow_eq} \\
\code{x == null} (ou \code{undefined}) & \code{narrow_keep_nullish_only} & \code{narrow_drop_null} (\code{_undef}) \\
\code{x} & \code{narrow_truthy} & \code{narrow_falsy} \\
\code{!x} & \code{narrow_falsy} & \code{narrow_truthy} \\
\hline
\end{tabular}
\caption{Motifs reconnus par \code{narrow_env_for_branch}
(\src{src/engine/cfg_analyzer.rs}{212}{293}) ; \code{x} est une variable,
\code{v} un littéral numérique. L'IR confond \code{==} et \code{===}.}
\label{tab:treillis-simples-motifs}
\end{table}

Tout autre motif laisse l'environnement inchangé, ce qui est sûr : littéral à
gauche (\code{10 > count}), comparaison entre deux variables
(\code{count < limit}, le faux positif du \cref{lst:interpretation-abstraite-limit}),
comparaison à une chaîne (\code{mode === "a"}), \code{typeof x === "string"},
conditions composées par \code{&&} ou \code{||} (qui sont des diamants de CFG,
chaque membre étant alors une garde séparée).

\begin{piege}
  La branche fausse de \code{x < 5} applique \code{narrow_geq(5)}, la
  négation \emph{numérique}. Or \code{undefined < 5} est faux : la branche
  \og sinon\fg{} est prise par une valeur qui n'est pas $\ge 5$. C'est
  pourtant sûr, parce que \code{StateValue} ne délègue la comparaison qu'à
  sa composante \code{num} et garde toutes les autres telles quelles
  (\src{src/domains/impls/state_value.rs}{604}{629}, avec le commentaire
  \og \code{null < 5} est vrai en JS\fg{}). Un contributeur qui ferait
  \og mourir\fg{} les sortes non numériques dans une branche de comparaison
  introduirait un faux négatif.
\end{piege}

\subsection{La véracité, et une garde impossible qui n'élague rien}

\code{if (n)} ne compare à rien ; la branche vraie exclut les valeurs
\emph{falsy} de JavaScript : \code{null}, \code{undefined}, \code{0},
\code{""}, \code{false} (et \code{NaN}, qu'aucun intervalle ne contient).
Côté nombres, \code{narrow_truthy} délègue à \code{narrow_neq(0.0)} :

\rustsnippet[label={lst:treillis-simples-truthy}]{src/domains/impls/state_value.rs}{565}{582}{le raffinement de véracité (branche vraie)}

Sur un point $[0,0]$, le résultat est $\Bot$ ; sur $[0,5]$, rien ne change.
Les composantes booléenne et chaîne sont traitées de même (\code{False}
meurt, \code{""} est retiré de l'ensemble). Deux programmes le montrent :

\begin{lstlisting}[language=TypeScript,caption={Garde de véracité (\file{/tmp/ch09/limits.tsx}, extrait)},label={lst:treillis-simples-truthy-ex}]
export function TruthyFromZero() {
  const [n, setN] = useState(0);
  useEffect(() => {
    if (n) setN(n + 1);
  }, [n]);
  return <div>{n}</div>;
}

export function TruthyFromOne() {
  const [n, setN] = useState(1);
  useEffect(() => {
    if (n) setN(n + 1);
  }, [n]);
  return <div>{n}</div>;
}
\end{lstlisting}

Dans \code{TruthyFromZero}, la branche raffine $[0,0]$ en $\Bot$, puis
\code{n + 1} vaut $\Bot$ (l'addition absorbe $\Bot$), et l'écriture ne change
rien : pas de boucle, conformément à la réalité (test
\code{truthy_guarded_increment_from_zero_never_starts}). Dans
\code{TruthyFromOne}, $[1,1]$ survit, et le compteur croît :

\begin{sortie}
$ reactant check /tmp/ch09/limits.tsx
  TruthyFromOne  (2 hooks)  /tmp/ch09/limits.tsx
    warn   infinite-loop  [hook:0]  (line 29:2)  this effect keeps pushing state `n` (its deps do not provably gate it, so the effect can re-run every render) to new values on every run. Potential infinite render loop
       (2 trace step(s), rerun with --trace)
  TruthyFromZero  (2 hooks)  /tmp/ch09/limits.tsx
    warn   redundant-set-state  [hook:0]  (line 21:2)  state `n` is set to a stable value it already holds, so the update is redundant
       (1 trace step(s), rerun with --trace)
\end{sortie}

(Sortie élaguée aux deux composants concernés.) Le \Warning{}
\regle{redundant-set-state} sur \code{TruthyFromZero} est un faux positif :
l'écriture n'a jamais lieu. La règle réévalue l'argument \code{n + 1} dans
l'environnement convergé \emph{sans} le raffinement de la branche, obtient
$[1,1]$, et conclut \og même valeur\fg{} parce que $[1,1]$ et $[0,0]$ sont tous
deux des valeurs \og stables\fg{} (\cref{chap:regles-etat}).

Enfin, un raffinement qui donne $\Bot$ ne rend pas la branche morte :
\code{narrow_env_for_branch} remplace la variable testée dans
l'environnement, mais le bloc reste analysé. Si l'écriture ne dépend pas de
la variable raffinée, elle a lieu. On le voit avec une première version du
\cref{lst:treillis-simples-shift} où l'effet s'écrit
\texttt{const y = x <{}< 1; if (y === 2) setCount(count + 1);} : \code{y} vaut
$[3,3]$, la garde le raffine en $\Bot$, et pourtant l'écriture
\code{count + 1} est prise en compte (\regle{infinite-loop} tire, fichier
\file{/tmp/ch09/shift.tsx}). Ne pas élaguer est sûr ; c'est seulement moins
précis qu'un élagage de bloc.

% ===========================================================================
\section{Bilan : la soundness opération par opération}
\label{sec:treillis-simples-soundness}

Le \cref{tab:treillis-simples-bilan} rassemble les arguments du chapitre.
Une opération \emph{exacte} ne perd aucune information
($\gam(a \join b) = \gam(a) \cup \gam(b)$, ou l'analogue) ; une opération
\emph{sound} sur-approxime ; un \emph{défaut} sous-approxime dans certains
cas, ce qui peut produire un faux négatif.

\begin{table}[htbp]\centering\small
\begin{tabular}{p{4.1cm}p{5.4cm}p{4.2cm}}
\hline
Opération & Argument & Statut \\
\hline
\code{BoolVal} : join, meet & bijection avec $\powerset(\Bool)$ & exact \\
\code{SetterVal} : join, meet & \cref{prop:treillis-simples-flat-treillis} & sound \\
\code{StrConst} : join, meet & \cref{prop:treillis-simples-strconst} & sound (join exact sous le seuil) \\
concaténation de chaînes & produit cartésien, puis \code{from_set} & exact sous le seuil \\
\code{hull}, \code{meet} des intervalles & \cref{prop:treillis-simples-interval-treillis} & sound \\
\code{add}, \code{sub}, \code{neg} & \cref{thm:treillis-simples-arith} & sound (bornes infinies) \\
\code{mul} & \cref{thm:treillis-simples-arith}, sauf $0 \times \pm\infty$ & \textbf{défaut} \\
\code{rem}, \code{pow} & \cref{prop:treillis-simples-rem}, coins du quadrant & sound (valeurs finies) \\
\code{/}, \code{&&}, \code{||} & $\Top$ & sound \\
comparaisons, \code{in}, \code{instanceof} & \code{BoolVal::Top} & sound \\
\code{&}, \code{|}, \code{^}, \texttt{>{}>{}>}, \code{~} & plage int32/uint32, bits & sound \\
\texttt{<{}<}, \texttt{>{}>} & multiplication / division des bornes brutes & \textbf{défaut} (non-entiers) \\
\code{widen}, \code{widen_to} & \cref{thm:treillis-simples-widen-to} & widening \\
\code{narrow_lt} \dots{} \code{narrow_neq} & \cref{thm:treillis-simples-narrow} & sound \\
\hline
\end{tabular}
\caption{Statut de soundness des opérations des treillis simples.}
\label{tab:treillis-simples-bilan}
\end{table}

Deux défauts sur quatorze lignes : c'est peu, et chacun tient à une
hypothèse implicite (pas de \code{NaN} dans un repli \code{min}/\code{max},
des bornes entières dans un décalage). Tous deux se corrigent dans la couche
intervalle, une fois pour toutes les règles, ce qui est l'esprit du
projet (principes de \file{CLAUDE.md} : corriger le moteur, pas les règles).

% ===========================================================================
\section{Limites et défauts observés}
\label{sec:treillis-simples-limites}

\subsection{$0 \times \pm\infty$ : un $\Bot$ inattendu}

Le repli de \code{mul} (\src{src/domains/impls/interval.rs}{180}{181})
calcule le minimum et le maximum des quatre produits de coins avec
\code{f64::min} et \code{f64::max}. Or ces fonctions \emph{ignorent}
\code{NaN} : \code{f64::min(x, NaN)} vaut \code{x}. Un produit de coins
$0 \times \pm\infty$, qui vaut \code{NaN}, est donc silencieusement écarté.
C'est souvent sans conséquence ($[0,5] \times [0,+\infty]$ donne bien
$[0,+\infty]$, les autres coins suffisent), mais si les \emph{quatre} coins
sont \code{NaN}, les replis rendent leurs valeurs initiales, $+\infty$ et
$-\infty$ : l'intervalle $[+\infty, -\infty]$, c'est-à-dire $\Bot$. Cela
arrive pour $[0,0] \times [-\infty, +\infty]$. Or, pour tout $x$ fini,
$0 \times x = 0$ : la valeur concrète $0$ est perdue.

\begin{lstlisting}[language=TypeScript,caption={Un produit par zéro d'un état non borné (\file{/tmp/ch09/limits.tsx}, extrait)},label={lst:treillis-simples-zero}]
export function ZeroTimesUnbounded() {
  const [n, setN] = useState(0);
  const [k, setK] = useState(0);
  const z = n * 0;
  useEffect(() => {
    setK(k + 1 + z);
  }, [k]);
  return (
    <div>
      <button onClick={() => setN(n + 1)}>+</button>
      <button onClick={() => setN(n - 1)}>-</button>
      {k}
    </div>
  );
}
\end{lstlisting}

Les deux handlers élargissent \code{n} dans les deux sens, vers
$[-\infty, +\infty]$, et \code{z = n * 0} est alors évalué à $\Bot$.
L'analyseur le laisse voir par un diagnostic contradictoire :

\begin{sortie}
$ reactant check /tmp/ch09/limits.tsx
  ZeroTimesUnbounded  (5 hooks)  /tmp/ch09/limits.tsx
    warn   infinite-loop  [hook:1]  (line 7:2)  this effect keeps pushing state `k` (its deps do not provably gate it, so the effect can re-run every render) to new values on every run. Potential infinite render loop
       (2 trace step(s), rerun with --trace)
    warn   missing-deps  [hook:2]  var:z  (line 7:2)  `z` is used in this effect but not in its deps array, and its value never changes between renders
       (1 trace step(s), rerun with --trace)
\end{sortie}

\og its value never changes\fg{} est la formulation commune aux valeurs de
stabilité $\Bot$ et \code{Stable} ; la règle, elle, ne se tait que sur
\code{Stable}. Rejouée avec \code{const z = 0;} au lieu de
\code{const z = n * 0;}, la même sortie ne porte plus que
l'\regle{infinite-loop} : c'est bien le $\Bot$ de \code{z} qui produit ce
\regle{missing-deps}. Ici la conséquence est un faux positif, et la boucle
sur \code{k} reste détectée. La raison n'est pas que \code{n} resterait
borné : les deux gestionnaires font grandir son intervalle des deux côtés à
chaque tour, et il atteint bien $[-\infty,+\infty]$ (c'est ce qui rend
\code{z} inatteignable). Mais le \code{widen_trace} qui alimente
\regle{widening-info} n'enregistre que les élargissements de
\code{state_from_render} et \code{state_from_effects}
(\src{src/engine/fixpoint.rs}{514}{516}, commentaire
\og handler widening is not a bug\fg{}) : l'élargissement de \code{n}, dû
aux gestionnaires de clic, n'y apparaît pas, quand bien même le \code{n}
final est bien non borné. Mais $\Bot$
signifie \og inatteignable\fg{} : toute écriture dont la valeur dépendrait
de \code{z} deviendrait un no-op abstrait, et une boucle réelle pourrait
disparaître. La correction naturelle (non appliquée au commit étudié) est
de traiter $0 \times \pm\infty$ comme $0$ (puisque $\pm\infty$ dénote
l'absence de borne et non une valeur), ou de rendre $\Top$ dès qu'un coin
est \code{NaN} ; c'est l'\cref{exo:treillis-simples-mul}.

\subsection{\code{Infinity} comme valeur}

La \cref{def:treillis-simples-intervalle} traite $\pm\infty$ comme absence
de borne. Mais \code{useState(1e999)} amorce le slot à
\code{point(f64::INFINITY)} $= [+\infty, +\infty]$, une vraie valeur. Suivons
\texttt{setN(n \% 3)} : le dividende est positif, le bit est faux (un infini
n'est pas entier), donc la borne est $\mu = 3$ et \code{rem} rend $[0, 3]$
(\cref{lst:treillis-simples-rem-pow}). En JavaScript,
\texttt{Infinity \% 3} vaut \code{NaN}, qui n'y est pas. Le dossier de notes a
confirmé ce calcul par une sonde, sans effet observable sur les diagnostics
(\code{NaN} est \code{Object.is}-égal à lui-même, la boucle s'arrête
concrètement). \code{docs/limitations.md} ne mentionne que les \code{NaN}
de \code{/}, de \texttt{\%} par un diviseur possiblement nul et de \code{**}
hors du quadrant positif. L'écart est une limite de la représentation : un
intervalle ne peut pas dire à la fois \og non borné\fg{} et \og contient
\code{Infinity}\fg{}.

\subsection{Les cycles finis}

Le domaine représente l'ensemble des valeurs \emph{atteignables}, pas les
\emph{transitions} entre valeurs. Un état qui cycle dans un ensemble fini a
donc un point fixe abstrait fini, atteint sans widening :

\begin{lstlisting}[language=TypeScript,caption={Deux cycles finis (\file{/tmp/ch09/limits.tsx}, extrait)},label={lst:treillis-simples-cycles}]
export function ModCycle() {
  const [i, setI] = useState(0);
  useEffect(() => {
    setI((i + 1) % 3);
  }, [i]);
  return <div>{i}</div>;
}

export function StrToggle() {
  const [mode, setMode] = useState("a");
  useEffect(() => {
    setMode(mode === "a" ? "b" : "a");
  }, [mode]);
  return <div>{mode}</div>;
}
\end{lstlisting}

\code{i} converge vers $[0,2]$ (\cref{ex:treillis-simples-mod}) et
\code{mode} vers $\{\code{"a"}, \code{"b"}\}$. Or les deux composants
bouclent concrètement : chaque écriture change la valeur (au sens
d'\code{Object.is}), donc la dep change, donc l'effet se réexécute.
L'analyseur les déclare sans défaut :

\begin{sortie}
$ reactant check /tmp/ch09/limits.tsx --show-clean
  ModCycle  (2 hooks)  /tmp/ch09/limits.tsx  ✓
  StrToggle  (2 hooks)  /tmp/ch09/limits.tsx  ✓
  …
\end{sortie}

(Sortie élaguée aux deux composants concernés.) Le bras \og point fixe\fg{}
de \regle{infinite-loop} exige que le slot ait été élargi et soit non borné,
ce qui n'est pas le cas d'un cycle fini ; le bras \og churn\fg{} ne
concerne que les références (\cref{chap:infinite-loop,chap:churn}). Aucun
treillis de ce chapitre ne peut porter l'information manquante, qui est une
propriété de transition (\og l'écriture diffère toujours de la valeur
lue\fg{}) ; c'est une limite de conception. Au commit \snapshotCommit{},
elle n'est documentée ni dans \file{docs/limitations.md}, ni par une issue
ouverte (aucune ne correspond à une recherche sur \og cycle\fg{}).

\subsection{Imprécisions assumées}

Pour mémoire, les imprécisions sûres rencontrées dans le chapitre :
comparaisons jamais évaluées (\code{1 < 2} vaut \code{BoolVal::Top}) ;
\code{narrow_neq} et \code{narrow_truthy} incapables de couper un intervalle
($[0,5]$ reste $[0,5]$ sous \code{if (n)}) ; raffinement limité aux motifs
\og variable contre littéral\fg{} ; garde impossible qui n'élague pas le
bloc ; \code{1 + "a"} évalué à $\Top$ au lieu d'une chaîne ;
\code{!}$\Bot$ évalué à un booléen inconnu ; seuils récoltés sans
discernement. Chacune peut produire un faux positif, aucune un faux négatif.

% ===========================================================================
\begin{resume}
\begin{itemize}
  \item Un domaine est un type qui implémente \code{AbstractDomain}
    (\file{src/domains/mod.rs}) ; l'ordre est \code{PartialOrd}. Le contrat
    implicite (\cref{def:treillis-simples-contrat}) exige que join, widening
    et raffinements majorent ; \code{meet} n'est pas utilisé par l'analyse.
  \item \code{flat_lattice!} (\file{src/domains/impls/mod.rs}) engendre les
    treillis plats \code{BoolVal} (exact) et \code{SetterVal} (identité
    exacte ou $\Top$). \code{StrConst} (\file{str_const.rs}) est un ensemble
    d'au plus 4 chaînes, $\Top$ au-delà ; hauteur finie, \code{widen = join}.
  \item \code{Interval} (\file{interval.rs}) : bornes \code{f64}, bit
    \code{is_int} hors égalité, $\Bot$ multiple (tester \code{is_bottom}),
    \code{NaN} jamais représenté. Join = enveloppe convexe.
  \item Arithmétique : formules classiques, sound grâce à la monotonie de
    l'arrondi ; \code{/} donne $\Top$ ; \texttt{\%} et \code{**} (commit
    \code{548f922}) rendent \code{None} quand \code{NaN} est possible ;
    opérateurs bit à bit bornés par int32/uint32 (\code{eval_bitwise}).
  \item Widening à seuils (\adr{14}) : chaque borne saute au seuil le plus
    serré parmi les littéraux du composant ; terminaison en au plus
    $|T| + 1$ sauts par borne. Chronologie du compteur : $[0,0]$, $[0,1]$,
    $[0,2]$, puis saut à $10$ ou à $+\infty$ au tour 3.
  \item Raffinement sur garde : \code{narrow_*}, choisis par
    \code{narrow_env_for_branch} selon la forme \og variable contre
    littéral\fg{} ; le resserrement entier n'a lieu que si \code{is_int} est
    prouvé. Pas de phase descendante.
  \item Défauts : $0 \times \pm\infty$ donne $\Bot$ ; \texttt{<{}<} et \texttt{>{}>}
    sur des non-entiers sous-approximent (faux négatif observé) ;
    \code{Infinity} comme valeur ; cycles finis invisibles.
\end{itemize}
Le \cref{chap:statevalue-stabilite} assemble ces treillis en un produit sur
les sortes JavaScript et y ajoute le treillis de stabilité référentielle ;
le \cref{chap:transfert-stores} décrit l'évaluateur qui les fait vivre.
\end{resume}

\section*{Exercices}
\addcontentsline{toc}{section}{Exercices}

\begin{exercice}{Meet des treillis plats}{treillis-simples-flat-meet}
  (a) Montrer que le \code{meet} de \code{flat_lattice!} est la plus grande
  borne inférieure. (b) Calculer \code{True} $\meet$ \code{False} et
  vérifier la condition 6 du contrat pour \code{BoolVal}. (c) Pourquoi le
  bras \code{(a, b) if a == b} doit-il précéder les autres dans
  \code{join} ?
\end{exercice}
\begin{solution}
  (a) Égaux : $a$ est le plus grand minorant de $\{a\}$. L'un est $\Top$ :
  l'autre est minorant et tout minorant commun le minore. Sinon, deux
  intermédiaires distincts (ou un intermédiaire et $\Bot$) n'ont que $\Bot$
  comme minorant commun. (b) $\Bot$ ; $\gam(\code{True}) \cap
  \gam(\code{False}) = \emptyset = \gam(\Bot)$ : exact. (c) Sans lui,
  $\Top \join \Top$ ou $m \join m$ tomberait dans le dernier bras et rendrait
  $\Top$ : correct pour $\Top$, mais $m \join m = \Top$ serait une perte
  de précision gratuite (un booléen toujours \code{True} deviendrait
  inconnu au premier join).
\end{solution}

\begin{exercice}{Chaînes}{treillis-simples-str}
  Calculer (a) $\{\code{"a"},\code{"b"}\} \join \{\code{"c"}\} \join
  \{\code{"d"}\} \join \{\code{"e"}\}$ en associant à gauche ;
  (b) $\{\code{"a"},\code{"b"}\} \meet \{\code{"c"}\}$ ; (c) la valeur de
  \code{"x-" + s} pour $s = \{\code{"a"},\code{"b"}\}$, puis pour
  $s = \{\code{"a"},\code{"b"},\code{"c"}\}$ et le préfixe
  $\{\code{"x"},\code{"y"}\}$.
\end{exercice}
\begin{solution}
  (a) $\{a,b,c\}$, puis $\{a,b,c,d\}$ (quatre éléments, encore
  représentable), puis $\Top$ (cinq). (b) Intersection vide :
  \code{from_set} rend $\Bot$. (c) $\{\code{"x-a"}, \code{"x-b"}\}$ ; le
  produit $\{x,y\} \times \{a,b,c\}$ a six éléments, \code{str_set} le
  remplace par \og une chaîne\fg{} ($\Top$ de \code{StrConst}).
\end{solution}

\begin{exercice}{Widening à seuils}{treillis-simples-widen}
  Avec $T = \{10, 20, 100\}$, calculer $[0,5] \widen_T [0,12]$, puis
  $[0,20] \widen_T [0,25]$, puis $[0,100] \widen_T [0,101]$, puis
  $[0,5] \widen_T [-3,5]$. Combien de sauts la borne haute peut-elle faire au
  plus ?
\end{exercice}
\begin{solution}
  $[0,20]$ (plus petit seuil $\ge 12$) ; $[0,100]$ ; $[0,+\infty]$ (aucun
  seuil $\ge 101$) ; $[-\infty, 5]$ (aucun seuil $\le -3$). La borne haute
  prend ses valeurs dans $\{10, 20, 100, +\infty\}$ après le premier saut :
  au plus $|T| + 1 = 4$ sauts.
\end{solution}

\begin{exercice}{Reste et puissance}{treillis-simples-rem-pow}
  Calculer avec le code : $[-7,8] \mathbin{\%} [3,3]$ ;
  $[0,9] \mathbin{\%} [-4,-2]$ ; $[1,2] \mathbin{\%} [5,5]$ ;
  $[0,9] \mathbin{\%} [0,3]$ ; $[0,2] \mathbin{**} [0,3]$ ;
  $[-1,2] \mathbin{**} [2,2]$. Commenter la précision du troisième.
\end{exercice}
\begin{solution}
  $[-2,2]$ ($\mu = 3$, $\beta = 2$) ; $[0,3]$ ($\mu = 4$, et en effet
  $7 \mathbin{\%} -4 = 3$) ; $[0,2]$ ; \code{None}, donc $\Top$ (diviseur
  possiblement nul) ; coins $0^0 = 1$, $0^3 = 0$, $2^0 = 1$, $2^3 = 8$,
  donc $[0,8]$ ; \code{None} (base possiblement négative), donc $\Top$. Le
  troisième est sound mais imprécis : le résultat exact est $[1,2]$, mais le
  code ne sait pas que $x \mathbin{\%} 5 = x$ quand $0 \le x < 5$ et met la
  borne basse à $0$.
\end{solution}

\begin{exercice}{Raffinements entiers}{treillis-simples-narrow}
  Sur $I = [0,9]$ entier, calculer les raffinements sous \code{x >= 2.5},
  \code{x > 7}, \code{x == 12}, et sous \code{x !== 9} pour $I = [9,9]$.
  Refaire le premier avec \code{is_int} faux. Prouver que
  \code{narrow_gt} est sound avec \code{is_int}.
\end{exercice}
\begin{solution}
  $[3,9]$ ($\lceil 2{,}5 \rceil$) ; $[8,9]$ ($\lfloor 7 \rfloor + 1$) ;
  $\Bot$ ($12 \notin [0,9]$) ; $\Bot$ (point égal à $9$). Sans bit :
  $[2{,}5; 9]$. Preuve : un entier $x > v$ vérifie
  $x \ge \lfloor v \rfloor + 1$, car $\lfloor v \rfloor + 1$ est le plus petit
  entier strictement supérieur à $v$ ; la nouvelle borne basse
  $\max(\ell, \lfloor v \rfloor + 1)$ est donc $\le x$.
\end{solution}

\begin{exercice}{Bit à bit}{treillis-simples-bit}
  Calculer avec \code{eval_bitwise} : $[0,300] \mathbin{\&} 255$ ;
  $[0,5] \mathbin{\hat{}} [0,2]$ ; $[-4,4] \mathbin{|} [0,1]$ ;
  $[100, 2^{30}] \ll 2$ ; $x \ggg 0$ pour $x$ quelconque.
\end{exercice}
\begin{solution}
  $[0,255]$. $\max = 5$, $\lfloor\log_2 5\rfloor + 1 = 3$ bits : $[0,7]$.
  Opérande gauche possiblement négatif : pas de raffinement, plage int32.
  $2^{30} \times 4 = 2^{32}$ déborde d'int32 : plage int32. $x \ggg 0 \in
  [0, 2^{32} - 1]$ (la plage uint32 est le raffinement même, $k = 0$).
\end{solution}

\begin{exercice}{Corriger la multiplication}{treillis-simples-mul}
  Proposer une correction de \code{Interval::mul} pour $0 \times \pm\infty$
  et prouver sa soundness sous l'hypothèse \og $\pm\infty$ est une absence de
  borne\fg{}.
\end{exercice}
\begin{solution}
  Calculer chaque produit de coins avec la convention $0 \times \pm\infty =
  0$ (par exemple \code{if p.is_nan() { 0.0 } else { p }}). Preuve : soit
  $x \in [a,b]$, $y \in [c,d]$ finis. Remplaçons une borne infinie par une
  borne finie $M$ assez grande pour contenir $x$ et $y$ : sur le pavé fini,
  le \cref{thm:treillis-simples-arith} s'applique, et les coins valant
  $0 \times M = 0$ sont exactement ceux que la convention remplace ; les coins
  $u \times M$ avec $u \neq 0$ ne sont pas plus serrés que $u \times \infty$.
  L'intervalle corrigé contient donc $xy$. Pour $[0,0] \times
  [-\infty,+\infty]$ il rend $[0,0]$, exact. L'alternative plus simple,
  rendre $\Top$ si un coin est \code{NaN}, est sound mais perd
  $[0,5] \times [0,+\infty] = [0,+\infty]$.
\end{solution}

\begin{exercice}{Corriger les décalages}{treillis-simples-shift}
  Corriger les bras \code{Shl} et \code{Shr} d'\code{eval_bitwise} pour les
  opérandes non entiers, et vérifier sur \code{ShlFloat} et \code{ShrFloat}
  (\cref{lst:treillis-simples-shift}).
\end{exercice}
\begin{solution}
  Sur la plage int32, $\code{ToInt32}(x) = \mathrm{trunc}(x)$, fonction
  croissante : l'image de $[\ell, h]$ est $[\mathrm{trunc}(\ell),
  \mathrm{trunc}(h)]$. Il suffit donc de remplacer \code{a.lo} et
  \code{a.hi} par \code{a.lo.trunc()} et \code{a.hi.trunc()} avant de
  multiplier ou de diviser, le reste de la preuve étant celle du cas entier
  (multiplication par $2^k$ sans débordement, division entière par $2^k$
  arrondie vers le bas, toutes deux croissantes). \code{ShlFloat} :
  $\mathrm{trunc}(1{,}5) = 1$, $y = [2,2]$, l'effet écrit \code{count + 1} :
  la boucle est vue. \code{ShrFloat} : $\mathrm{trunc}(-1{,}5) = -1$,
  $y = [-1,-1]$, idem. La correction est centrale : toutes les règles qui
  lisent une valeur issue d'un décalage en profitent.
\end{solution}
